\chapter{Environment and configuration reference}\label{ch:ref-env}

TandemLLM reads four families of settings. Engine switches are \code{QWEN38_*} variables, read through one registry. Server settings are command-line flags (Chapter~\ref{ch:ref-cli}); \code{ops/start.sh}\index{ops/start.sh@\texttt{ops/start.sh}} builds them from \code{ops/serve.env}\index{ops/serve.env@\texttt{ops/serve.env}}, and a few \code{QSE_*} variables set flags' defaults or serve the scripts. The proxy, the collector and the plugins read \code{TANDEM_*} variables. And the conductor's behaviour lives in its versioned documents, which the dashboard edits.

\section{The served profile}\index{serve.env}

\code{ops/serve.env} is the served configuration. On an installed engine it is \code{~/TandemLLM/run/ops/serve.env}, rewritten by every run of the installer, and \code{local.env} beside it is read last for your own values.

% Generated by manual/tools/gen_reference.py from the code. Do not edit by hand.

\begin{reftable}{The keys ops/serve.env sets, in file order (paths not shown)}{tab:ref-serve-env}{\hd{Key} & \hd{Value} & \hd{Kind}}
\texttt{HOME\_\allowbreak{}DIR} & \texttt{(a path)} & read by ops/start.sh \\
\texttt{REPO} & \texttt{(a path)} & read by ops/start.sh \\
\texttt{PY} & \texttt{(a path)} & read by ops/start.sh \\
\texttt{PYTHONPATH} & \texttt{(a path)} & read by ops/start.sh \\
\texttt{HOST} & \texttt{0.0.0.0} & read by ops/start.sh \\
\texttt{PORT} & \texttt{8000} & read by ops/start.sh \\
\texttt{SERVED\_\allowbreak{}MODEL} & \texttt{qwen38-\allowbreak{}spark-\allowbreak{}engine} & read by ops/start.sh \\
\texttt{NV} & \texttt{(a path)} & read by ops/start.sh \\
\texttt{HEAD} & \texttt{(a path)} & read by ops/start.sh \\
\texttt{CKPT8} & \texttt{(a path)} & read by ops/start.sh \\
\texttt{CKPT16} & \texttt{(a path)} & read by ops/start.sh \\
\texttt{CORPUS} & \texttt{(a path)} & read by ops/start.sh \\
\texttt{MAX\_\allowbreak{}LEN} & \texttt{262144} & read by ops/start.sh \\
\texttt{DEFAULT\_\allowbreak{}MAX\_\allowbreak{}TOKENS} & \texttt{32768} & read by ops/start.sh \\
\texttt{PATTERN\_\allowbreak{}STOP} & \texttt{(empty)} & read by ops/start.sh \\
\texttt{CONDUCTOR} & \texttt{on} & read by ops/start.sh \\
\texttt{REASONING\_\allowbreak{}FORMAT} & \texttt{tags} & read by ops/start.sh \\
\texttt{REASONING\_\allowbreak{}EFFORT} & \texttt{medium} & read by ops/start.sh \\
\texttt{LEN\_\allowbreak{}FIXED} & \texttt{0} & read by ops/start.sh \\
\texttt{LEN\_\allowbreak{}LATCH} & \texttt{1} & read by ops/start.sh \\
\texttt{BUDGET} & \texttt{16} & read by ops/start.sh \\
\texttt{MEM\_\allowbreak{}CEILING\_\allowbreak{}GIB} & \texttt{70} & read by ops/start.sh \\
\texttt{MEM\_\allowbreak{}PEAK\_\allowbreak{}GIB} & \texttt{77} & read by ops/start.sh \\
\texttt{RESERVE\_\allowbreak{}SLACK\_\allowbreak{}GIB} & \texttt{2} & read by ops/start.sh \\
\texttt{RESIDENT\_\allowbreak{}GB} & \texttt{2.5} & read by ops/start.sh \\
\texttt{RESIDENT\_\allowbreak{}STASH\_\allowbreak{}GB} & \texttt{1.0} & read by ops/start.sh \\
\texttt{RESIDENT\_\allowbreak{}GUEST\_\allowbreak{}MAX} & \texttt{8192} & read by ops/start.sh \\
\texttt{RESIDENT\_\allowbreak{}PIN\_\allowbreak{}GB} & \texttt{4.0} & read by ops/start.sh \\
\texttt{RESIDENT\_\allowbreak{}PARK\_\allowbreak{}GB} & \texttt{4} & read by ops/start.sh \\
\texttt{CACHE\_\allowbreak{}GB} & \texttt{2} & read by ops/start.sh \\
\texttt{STORE\_\allowbreak{}ENTRY\_\allowbreak{}GB} & \texttt{1} & read by ops/start.sh \\
\texttt{NARROW\_\allowbreak{}ARM} & \texttt{1} & read by ops/start.sh \\
\texttt{END\_\allowbreak{}ANCHOR} & \texttt{on} & read by ops/start.sh \\
\texttt{END\_\allowbreak{}ALIGN} & \texttt{on} & read by ops/start.sh \\
\texttt{DROP\_\allowbreak{}PAGE\_\allowbreak{}CACHE} & \texttt{1} & read by ops/start.sh \\
\texttt{REQUEST\_\allowbreak{}TIMEOUT} & \texttt{3600} & read by ops/start.sh \\
\texttt{MAX\_\allowbreak{}QUEUE} & \texttt{8} & read by ops/start.sh \\
\texttt{QUEUE\_\allowbreak{}TIMEOUT} & \texttt{240} & read by ops/start.sh \\
\texttt{FUSE\_\allowbreak{}PROJ} & \texttt{0} & read by ops/start.sh \\
\texttt{QWEN38\_\allowbreak{}DF2\_\allowbreak{}TREE\_\allowbreak{}MODE} & \texttt{nodes} & engine switch \\
\texttt{QWEN38\_\allowbreak{}TREE\_\allowbreak{}ALIAS\_\allowbreak{}STATE} & \texttt{0} & engine switch \\
\texttt{DROP\_\allowbreak{}IDLE} & \texttt{1} & read by ops/start.sh \\
\texttt{QWEN38\_\allowbreak{}DRAFT\_\allowbreak{}TEMP} & \texttt{0.3} & engine switch \\
\texttt{SAMPLED\_\allowbreak{}TREE} & \texttt{det} & read by ops/start.sh \\
\texttt{QWEN38\_\allowbreak{}FUSED\_\allowbreak{}GDNPREFILL} & \texttt{1} & engine switch \\
\texttt{QWEN38\_\allowbreak{}DECODE\_\allowbreak{}ATTN} & \texttt{1} & engine switch \\
\texttt{QWEN38\_\allowbreak{}DRAFT\_\allowbreak{}NVFP4} & \texttt{1} & engine switch \\
\texttt{QWEN38\_\allowbreak{}NVFP4\_\allowbreak{}PREFILL\_\allowbreak{}V2} & \texttt{1} & engine switch \\
\texttt{QWEN38\_\allowbreak{}NVFP4\_\allowbreak{}PREFILL\_\allowbreak{}V2\_\allowbreak{}UNTIL} & \texttt{1025} & engine switch \\
\texttt{QWEN38\_\allowbreak{}NVFP4\_\allowbreak{}SKINNY} & \texttt{1} & engine switch \\
\texttt{QWEN38\_\allowbreak{}SKINNY\_\allowbreak{}TILES} & \texttt{(a path)} & engine switch \\
\texttt{QWEN38\_\allowbreak{}FUSED\_\allowbreak{}COMMIT} & \texttt{1} & engine switch \\
\texttt{QWEN38\_\allowbreak{}FUSED\_\allowbreak{}GDNVERIFY} & \texttt{1} & engine switch \\
\texttt{QWEN38\_\allowbreak{}DRAFT\_\allowbreak{}FC\_\allowbreak{}NVFP4} & \texttt{1} & engine switch \\
\texttt{QWEN38\_\allowbreak{}FUSED\_\allowbreak{}ADDNORM} & \texttt{1} & engine switch \\
\texttt{QWEN38\_\allowbreak{}TREE\_\allowbreak{}HOST\_\allowbreak{}DEPTH} & \texttt{1} & engine switch \\
\texttt{QWEN38\_\allowbreak{}HOST\_\allowbreak{}WALK} & \texttt{1} & engine switch \\
\texttt{QWEN38\_\allowbreak{}VERIFY\_\allowbreak{}GRAPH} & \texttt{1} & engine switch \\
\texttt{QWEN38\_\allowbreak{}COMMIT\_\allowbreak{}IN\_\allowbreak{}VERIFY} & \texttt{1} & engine switch \\
\texttt{QWEN38\_\allowbreak{}DRAFT\_\allowbreak{}GRAPH} & \texttt{1} & engine switch \\
\texttt{QWEN38\_\allowbreak{}GDNV\_\allowbreak{}ONE\_\allowbreak{}WARP} & \texttt{1} & engine switch \\
\texttt{QWEN38\_\allowbreak{}TREE\_\allowbreak{}NODES\_\allowbreak{}NARROW} & \texttt{16} & engine switch \\
\texttt{QWEN38\_\allowbreak{}TREE\_\allowbreak{}MS} & \texttt{8:\allowbreak{}77.9,\allowbreak{}16:\allowbreak{}78.9,\allowbreak{}24:\allowbreak{}87.0,\allowbreak{}32:\allowbreak{}97.0} & engine switch \\
\texttt{QWEN38\_\allowbreak{}VERIFY\_\allowbreak{}ROWS} & \texttt{32} & engine switch \\
\texttt{QWEN38\_\allowbreak{}SKINNY\_\allowbreak{}TILES\_\allowbreak{}WIDE} & \texttt{(a path)} & engine switch \\
\texttt{QWEN38\_\allowbreak{}TREE\_\allowbreak{}NODES} & \texttt{24} & engine switch \\
\texttt{QWEN38\_\allowbreak{}TREE\_\allowbreak{}WIDE\_\allowbreak{}AFTER} & \texttt{32} & engine switch \\
\texttt{QWEN38\_\allowbreak{}DEEP} & \texttt{32} & engine switch \\
\texttt{QWEN38\_\allowbreak{}DEEP\_\allowbreak{}AFTER} & \texttt{2} & engine switch \\
\texttt{QWEN38\_\allowbreak{}LEN\_\allowbreak{}SWITCH} & \texttt{1} & engine switch \\
\texttt{QWEN38\_\allowbreak{}LEN\_\allowbreak{}MODE} & \texttt{wide} & engine switch \\
\texttt{QWEN38\_\allowbreak{}STAIR\_\allowbreak{}TABLES} & \texttt{(a path)} & engine switch \\
\texttt{QWEN38\_\allowbreak{}STAIR\_\allowbreak{}RHO} & \texttt{(a path)} & engine switch \\
\texttt{QWEN38\_\allowbreak{}STAIR\_\allowbreak{}SKIP} & \texttt{1} & engine switch \\
\texttt{QWEN38\_\allowbreak{}SKINNY\_\allowbreak{}SRUN} & \texttt{1} & engine switch \\
\texttt{QSE\_\allowbreak{}USAGE\_\allowbreak{}LEDGER} & \texttt{(a path)} & server setting \\
\texttt{QSE\_\allowbreak{}CONDUCTOR\_\allowbreak{}CONFIG} & \texttt{(a path)} & server setting \\
\texttt{QSE\_\allowbreak{}BACKENDS} & \texttt{(empty)} & server setting \\
\texttt{QSE\_\allowbreak{}PARALLEL} & \texttt{1} & server setting \\
\texttt{QSE\_\allowbreak{}SIDE\_\allowbreak{}MAX\_\allowbreak{}LEN} & \texttt{65536} & server setting \\
\texttt{QSE\_\allowbreak{}PARALLEL\_\allowbreak{}FLOOR\_\allowbreak{}GB} & \texttt{15} & server setting \\
\texttt{QWEN38\_\allowbreak{}GDNV\_\allowbreak{}CONV\_\allowbreak{}BLOCK} & \texttt{32} & engine switch \\
\texttt{QWEN38\_\allowbreak{}GDNV\_\allowbreak{}CONV\_\allowbreak{}WARPS} & \texttt{1} & engine switch \\
\texttt{QWEN38\_\allowbreak{}GDNV\_\allowbreak{}WY} & \texttt{1} & engine switch \\
\texttt{QWEN38\_\allowbreak{}GDNV\_\allowbreak{}WY\_\allowbreak{}FUSED} & \texttt{1} & engine switch \\
\texttt{QWEN38\_\allowbreak{}GDNV\_\allowbreak{}WY\_\allowbreak{}MAXT} & \texttt{16} & engine switch \\
\texttt{QWEN38\_\allowbreak{}GDNV\_\allowbreak{}WY\_\allowbreak{}CHAIN\_\allowbreak{}MAXT} & \texttt{16} & engine switch \\
\texttt{QWEN38\_\allowbreak{}GDNV\_\allowbreak{}WY\_\allowbreak{}FUSED\_\allowbreak{}MAXT} & \texttt{16} & engine switch \\
\texttt{LIVE\_\allowbreak{}ACTIVITY} & \texttt{on} & read by ops/start.sh \\
\texttt{QSE\_\allowbreak{}DASHBOARD\_\allowbreak{}LOGIN} & \texttt{off} & server setting \\
\texttt{ROLES\_\allowbreak{}AUTOLOAD} & \texttt{(empty)} & read by ops/start.sh \\
\texttt{QWEN38\_\allowbreak{}DECIDER\_\allowbreak{}MODELS} & \texttt{(empty)} & engine switch \\
\texttt{QWEN38\_\allowbreak{}DECIDER\_\allowbreak{}DIR} & \texttt{(a path)} & engine switch \\
\texttt{QWEN38\_\allowbreak{}NARRATOR\_\allowbreak{}MODEL} & \texttt{(empty)} & engine switch \\
\texttt{QWEN38\_\allowbreak{}NARRATOR\_\allowbreak{}DIR} & \texttt{(empty)} & engine switch \\
\texttt{QWEN38\_\allowbreak{}SHADOW\_\allowbreak{}MODEL} & \texttt{(empty)} & engine switch \\
\texttt{QWEN38\_\allowbreak{}SHADOW\_\allowbreak{}DIR} & \texttt{(empty)} & engine switch \\
\texttt{QWEN38\_\allowbreak{}SHADOW\_\allowbreak{}GB} & \texttt{4} & engine switch \\
\texttt{QWEN38\_\allowbreak{}SHADOW\_\allowbreak{}GRID} & \texttt{512} & engine switch \\
\texttt{QWEN38\_\allowbreak{}SHADOW\_\allowbreak{}DISK} & \texttt{auto} & engine switch \\
\texttt{QWEN38\_\allowbreak{}SHADOW\_\allowbreak{}DISK\_\allowbreak{}GB} & \texttt{20} & engine switch \\
\texttt{QWEN38\_\allowbreak{}PROBE\_\allowbreak{}SNAPSHOTS} & \texttt{off} & engine switch \\
\texttt{QWEN38\_\allowbreak{}PROBE\_\allowbreak{}DIR} & \texttt{(a path)} & engine switch \\
\texttt{QSE\_\allowbreak{}DISK\_\allowbreak{}CACHE} & \texttt{(a path)} & server setting \\
\texttt{SESSION\_\allowbreak{}SPILL} & \texttt{on} & read by ops/start.sh \\
\texttt{SESSION\_\allowbreak{}SPILL\_\allowbreak{}GB} & \texttt{64} & read by ops/start.sh \\
\texttt{SESSION\_\allowbreak{}SPILL\_\allowbreak{}CODEC} & \texttt{zstd} & read by ops/start.sh \\
\texttt{DISK\_\allowbreak{}WRITE\_\allowbreak{}BUDGET\_\allowbreak{}GB} & \texttt{200} & read by ops/start.sh \\
\texttt{VISION} & \texttt{on} & read by ops/start.sh \\
\texttt{RESPONSE\_\allowbreak{}CACHE} & \texttt{off} & read by ops/start.sh \\
\texttt{QSE\_\allowbreak{}STATUS\_\allowbreak{}PORT} & \texttt{(empty)} & server setting \\
\texttt{QSE\_\allowbreak{}PROFILE\_\allowbreak{}NAME} & \texttt{(a path)} & server setting \\
\texttt{QSE\_\allowbreak{}SETUP\_\allowbreak{}ENV\_\allowbreak{}FILE} & \texttt{(a path)} & server setting \\
\end{reftable}


\section{Engine switches}\index{QWEN38 variables}

\code{engine/settings.py} lists every engine switch with its default and the module that reads it; \code{tests/test_settings.py} fails when a module reads one another way. The table shows each switch, its code default, and the value the served profile sets when it sets one.

% Generated by manual/tools/gen_reference.py from the code. Do not edit by hand.

\begin{reftable}{Every QWEN38\_* switch in engine/settings.py, with the value ops/serve.env sets}{tab:ref-knobs}{\hd{Variable} & \hd{Code default} & \hd{Served value \textbullet{} read in}}
\texttt{QWEN38\_\allowbreak{}ATTN\_\allowbreak{}STAGES3} & \texttt{1} & \texttt{tools/\allowbreak{}attn\_\allowbreak{}kernels.py} \\
\texttt{QWEN38\_\allowbreak{}BLOCKING\_\allowbreak{}SYNC} & \texttt{0} & \texttt{server/\allowbreak{}app.py} \\
\texttt{QWEN38\_\allowbreak{}CLIENT\_\allowbreak{}CWD} & \texttt{note} & \texttt{server/\allowbreak{}conductor/\allowbreak{}router.py} \\
\texttt{QWEN38\_\allowbreak{}CLIENT\_\allowbreak{}REASONING} & \texttt{*=\allowbreak{}reasoning\_\allowbreak{}content,\allowbreak{}hermes=\allowbreak{}tags} & \texttt{server/\allowbreak{}compat.py} \\
\texttt{QWEN38\_\allowbreak{}COMMIT\_\allowbreak{}IN\_\allowbreak{}VERIFY} & \texttt{0} & \texttt{1} \textbullet{} \texttt{engine/\allowbreak{}model.py} \\
\texttt{QWEN38\_\allowbreak{}CONDUCTOR\_\allowbreak{}FIT} & \texttt{model} & \texttt{server/\allowbreak{}conductor/\allowbreak{}modelfit.py} \\
\texttt{QWEN38\_\allowbreak{}CORPUS} & \texttt{none} & \texttt{engine/\allowbreak{}drafters/\allowbreak{}ngram.py} \\
\texttt{QWEN38\_\allowbreak{}DECIDER\_\allowbreak{}COLLATE} & \texttt{call} & \texttt{server/\allowbreak{}decider\_\allowbreak{}api.py} \\
\texttt{QWEN38\_\allowbreak{}DECIDER\_\allowbreak{}DEVICE} & \texttt{auto} & \texttt{server/\allowbreak{}decider\_\allowbreak{}api.py} \\
\texttt{QWEN38\_\allowbreak{}DECIDER\_\allowbreak{}DIR} & \texttt{\textasciitilde{}/\allowbreak{}.qwen38-\allowbreak{}spark-\allowbreak{}engine/\allowbreak{}deciders/\allowbreak{}laya} & \texttt{(a path)} \textbullet{} \texttt{server/\allowbreak{}decider\_\allowbreak{}api.py} \\
\texttt{QWEN38\_\allowbreak{}DECIDER\_\allowbreak{}MAX\_\allowbreak{}BATCH} & \texttt{64} & \texttt{server/\allowbreak{}decider\_\allowbreak{}api.py} \\
\texttt{QWEN38\_\allowbreak{}DECIDER\_\allowbreak{}MAX\_\allowbreak{}QUEUE} & \texttt{256} & \texttt{server/\allowbreak{}decider\_\allowbreak{}api.py} \\
\texttt{QWEN38\_\allowbreak{}DECIDER\_\allowbreak{}MODELS} & \texttt{none} & \texttt{server/\allowbreak{}decider\_\allowbreak{}api.py} \\
\texttt{QWEN38\_\allowbreak{}DECIDER\_\allowbreak{}WAIT\_\allowbreak{}MS} & \texttt{2} & \texttt{server/\allowbreak{}decider\_\allowbreak{}api.py} \\
\texttt{QWEN38\_\allowbreak{}DECODE\_\allowbreak{}ATTN} & \texttt{0} & \texttt{1} \textbullet{} \texttt{engine/\allowbreak{}model.py} \\
\texttt{QWEN38\_\allowbreak{}DEEP} & \texttt{0} & \texttt{32} \textbullet{} \texttt{server/\allowbreak{}app.py} \\
\texttt{QWEN38\_\allowbreak{}DEEP\_\allowbreak{}AFTER} & \texttt{2} & \texttt{2} \textbullet{} \texttt{server/\allowbreak{}app.py} \\
\texttt{QWEN38\_\allowbreak{}DF2\_\allowbreak{}TEMP} & \texttt{1.0} & \texttt{engine/\allowbreak{}drafters/\allowbreak{}dflash2.py} \\
\texttt{QWEN38\_\allowbreak{}DF2\_\allowbreak{}TREE\_\allowbreak{}MODE} & \texttt{paths} & \texttt{nodes} \textbullet{} \texttt{engine/\allowbreak{}drafters/\allowbreak{}dflash2.py} \\
\texttt{QWEN38\_\allowbreak{}DFLASH2} & \texttt{none} & \texttt{engine/\allowbreak{}drafters/\allowbreak{}dflash2.py} \\
\texttt{QWEN38\_\allowbreak{}DRAFT\_\allowbreak{}EXL3} & \texttt{none} & \texttt{engine/\allowbreak{}drafters/\allowbreak{}dflash2.py} \\
\texttt{QWEN38\_\allowbreak{}DRAFT\_\allowbreak{}FC\_\allowbreak{}NVFP4} & \texttt{0} & \texttt{1} \textbullet{} \texttt{engine/\allowbreak{}drafters/\allowbreak{}dflash2.py} \\
\texttt{QWEN38\_\allowbreak{}DRAFT\_\allowbreak{}GRAPH} & \texttt{0} & \texttt{1} \textbullet{} \texttt{engine/\allowbreak{}drafters/\allowbreak{}dflash2.py} \\
\texttt{QWEN38\_\allowbreak{}DRAFT\_\allowbreak{}HEAD} & \texttt{none} & \texttt{engine/\allowbreak{}drafters/\allowbreak{}dflash2.py} \\
\texttt{QWEN38\_\allowbreak{}DRAFT\_\allowbreak{}HEAD\_\allowbreak{}EXL3} & \texttt{none} & \texttt{engine/\allowbreak{}drafters/\allowbreak{}dflash2.py} \\
\texttt{QWEN38\_\allowbreak{}DRAFT\_\allowbreak{}HEAD\_\allowbreak{}NVFP4} & \texttt{0} & \texttt{engine/\allowbreak{}drafters/\allowbreak{}dflash2.py} \\
\texttt{QWEN38\_\allowbreak{}DRAFT\_\allowbreak{}NVFP4} & \texttt{0} & \texttt{1} \textbullet{} \texttt{engine/\allowbreak{}drafters/\allowbreak{}dflash2.py} \\
\texttt{QWEN38\_\allowbreak{}DRAFT\_\allowbreak{}TEMP} & \texttt{1.0} & \texttt{0.3} \textbullet{} \texttt{engine/\allowbreak{}drafters/\allowbreak{}dflash2.py} \\
\texttt{QWEN38\_\allowbreak{}DSPARK} & \texttt{none} & \texttt{engine/\allowbreak{}drafters/\allowbreak{}dspark.py} \\
\texttt{QWEN38\_\allowbreak{}DSPARK\_\allowbreak{}NO\_\allowbreak{}YARN} & \texttt{none} & \texttt{engine/\allowbreak{}drafters/\allowbreak{}dspark.py} \\
\texttt{QWEN38\_\allowbreak{}END\_\allowbreak{}ALIGN\_\allowbreak{}MOVED} & \texttt{on} & \texttt{server/\allowbreak{}endalign.py} \\
\texttt{QWEN38\_\allowbreak{}END\_\allowbreak{}ALIGN\_\allowbreak{}REASONING} & \texttt{on} & \texttt{server/\allowbreak{}endalign.py} \\
\texttt{QWEN38\_\allowbreak{}EXL3} & \texttt{none} & \texttt{engine/\allowbreak{}loader.py} \\
\texttt{QWEN38\_\allowbreak{}EXL3\_\allowbreak{}CHUNK\_\allowbreak{}UNTIL} & \texttt{128} & \texttt{tools/\allowbreak{}exl3\_\allowbreak{}linear.py} \\
\texttt{QWEN38\_\allowbreak{}EXL3\_\allowbreak{}FUSE} & \texttt{1} & \texttt{engine/\allowbreak{}loader.py} \\
\texttt{QWEN38\_\allowbreak{}EXL3\_\allowbreak{}MLP\_\allowbreak{}BRIDGE} & \texttt{1} & \texttt{tools/\allowbreak{}exl3\_\allowbreak{}linear.py} \\
\texttt{QWEN38\_\allowbreak{}EXL3\_\allowbreak{}ROWGRID} & \texttt{1} & \texttt{tools/\allowbreak{}exl3\_\allowbreak{}linear.py} \\
\texttt{QWEN38\_\allowbreak{}EXL3\_\allowbreak{}SILU\_\allowbreak{}FUSED} & \texttt{0} & \texttt{tools/\allowbreak{}exl3\_\allowbreak{}linear.py} \\
\texttt{QWEN38\_\allowbreak{}EXL3\_\allowbreak{}STAGES} & \texttt{3} & \texttt{tools/\allowbreak{}exl3\_\allowbreak{}linear.py} \\
\texttt{QWEN38\_\allowbreak{}EXL3\_\allowbreak{}TILES} & \texttt{none} & \texttt{tools/\allowbreak{}exl3\_\allowbreak{}linear.py} \\
\texttt{QWEN38\_\allowbreak{}FN\_\allowbreak{}DENSE} & \texttt{0} & \texttt{engine/\allowbreak{}models/\allowbreak{}flashnext/\allowbreak{}engine.py} \\
\texttt{QWEN38\_\allowbreak{}FN\_\allowbreak{}DRAFT\_\allowbreak{}HOT} & \texttt{32768} & \texttt{engine/\allowbreak{}models/\allowbreak{}flashnext/\allowbreak{}serving.py} \\
\texttt{QWEN38\_\allowbreak{}FN\_\allowbreak{}DROP\_\allowbreak{}SHARD\_\allowbreak{}CACHE} & \texttt{1} & \texttt{engine/\allowbreak{}models/\allowbreak{}flashnext/\allowbreak{}load.py} \\
\texttt{QWEN38\_\allowbreak{}FN\_\allowbreak{}FOLD\_\allowbreak{}COMMIT} & \texttt{1} & \texttt{engine/\allowbreak{}models/\allowbreak{}flashnext/\allowbreak{}engine.py} \\
\texttt{QWEN38\_\allowbreak{}FN\_\allowbreak{}GEMV\_\allowbreak{}SPLIT} & \texttt{1} & \texttt{tools/\allowbreak{}fn\_\allowbreak{}kernels.py} \\
\texttt{QWEN38\_\allowbreak{}FN\_\allowbreak{}LOOKUP} & \texttt{0} & \texttt{engine/\allowbreak{}models/\allowbreak{}flashnext/\allowbreak{}serving.py} \\
\texttt{QWEN38\_\allowbreak{}FN\_\allowbreak{}MIX} & \texttt{int8} & \texttt{engine/\allowbreak{}models/\allowbreak{}flashnext/\allowbreak{}pack.py} \\
\texttt{QWEN38\_\allowbreak{}FN\_\allowbreak{}MOE\_\allowbreak{}ST} & \texttt{3} & \texttt{tools/\allowbreak{}fn\_\allowbreak{}kernels.py} \\
\texttt{QWEN38\_\allowbreak{}FN\_\allowbreak{}MOE\_\allowbreak{}WARPS} & \texttt{4} & \texttt{tools/\allowbreak{}fn\_\allowbreak{}kernels.py} \\
\texttt{QWEN38\_\allowbreak{}FN\_\allowbreak{}PDL} & \texttt{1} & \texttt{tools/\allowbreak{}fn\_\allowbreak{}kernels.py} \\
\texttt{QWEN38\_\allowbreak{}FN\_\allowbreak{}PREFILL} & \texttt{kernels} & \texttt{engine/\allowbreak{}models/\allowbreak{}flashnext/\allowbreak{}engine.py} \\
\texttt{QWEN38\_\allowbreak{}FN\_\allowbreak{}PREFILL\_\allowbreak{}CHUNK} & \texttt{4096} & \texttt{engine/\allowbreak{}models/\allowbreak{}flashnext/\allowbreak{}engine.py} \\
\texttt{QWEN38\_\allowbreak{}FN\_\allowbreak{}PREFILL\_\allowbreak{}MM} & \texttt{fp32} & \texttt{tools/\allowbreak{}fn\_\allowbreak{}prefill\_\allowbreak{}kernels.py} \\
\texttt{QWEN38\_\allowbreak{}FN\_\allowbreak{}TABLE} & \texttt{warm} & \texttt{engine/\allowbreak{}models/\allowbreak{}flashnext/\allowbreak{}load.py} \\
\texttt{QWEN38\_\allowbreak{}FN\_\allowbreak{}TABLE\_\allowbreak{}CACHE\_\allowbreak{}GB} & \texttt{2} & \texttt{engine/\allowbreak{}models/\allowbreak{}flashnext/\allowbreak{}load.py} \\
\texttt{QWEN38\_\allowbreak{}FN\_\allowbreak{}WARM\_\allowbreak{}TABLE} & \texttt{1} & \texttt{engine/\allowbreak{}models/\allowbreak{}flashnext/\allowbreak{}load.py} \\
\texttt{QWEN38\_\allowbreak{}FP8\_\allowbreak{}HEAD} & \texttt{none} & \texttt{engine/\allowbreak{}loader.py} \\
\texttt{QWEN38\_\allowbreak{}FP8\_\allowbreak{}TILES} & \texttt{none} & \texttt{tools/\allowbreak{}fp8\_\allowbreak{}linear.py} \\
\texttt{QWEN38\_\allowbreak{}FUSED\_\allowbreak{}ADDNORM} & \texttt{0} & \texttt{1} \textbullet{} \texttt{engine/\allowbreak{}model.py} \\
\texttt{QWEN38\_\allowbreak{}FUSED\_\allowbreak{}ATTN} & \texttt{0} & \texttt{engine/\allowbreak{}model.py} \\
\texttt{QWEN38\_\allowbreak{}FUSED\_\allowbreak{}ATTN\_\allowbreak{}PREP} & \texttt{0} & \texttt{engine/\allowbreak{}model.py} \\
\texttt{QWEN38\_\allowbreak{}FUSED\_\allowbreak{}COMMIT} & \texttt{0} & \texttt{1} \textbullet{} \texttt{engine/\allowbreak{}model.py} \\
\texttt{QWEN38\_\allowbreak{}FUSED\_\allowbreak{}GDN} & \texttt{1} & \texttt{engine/\allowbreak{}model.py} \\
\texttt{QWEN38\_\allowbreak{}FUSED\_\allowbreak{}GDNBLOCK} & \texttt{1} & \texttt{engine/\allowbreak{}model.py} \\
\texttt{QWEN38\_\allowbreak{}FUSED\_\allowbreak{}GDNPRE} & \texttt{1} & \texttt{engine/\allowbreak{}model.py} \\
\texttt{QWEN38\_\allowbreak{}FUSED\_\allowbreak{}GDNPREFILL} & \texttt{0} & \texttt{1} \textbullet{} \texttt{engine/\allowbreak{}model.py} \\
\texttt{QWEN38\_\allowbreak{}FUSED\_\allowbreak{}GDNTREE} & \texttt{1} & \texttt{engine/\allowbreak{}model.py} \\
\texttt{QWEN38\_\allowbreak{}FUSED\_\allowbreak{}GDNVERIFY} & \texttt{0} & \texttt{1} \textbullet{} \texttt{engine/\allowbreak{}model.py} \\
\texttt{QWEN38\_\allowbreak{}FUSED\_\allowbreak{}HEAD} & \texttt{1} & \texttt{engine/\allowbreak{}model.py} \\
\texttt{QWEN38\_\allowbreak{}FUSED\_\allowbreak{}NORM} & \texttt{1} & \texttt{engine/\allowbreak{}model.py} \\
\texttt{QWEN38\_\allowbreak{}FUSE\_\allowbreak{}PROJ} & \texttt{0} & \texttt{tools/\allowbreak{}nvfp4\_\allowbreak{}linear\_\allowbreak{}v2.py} \\
\texttt{QWEN38\_\allowbreak{}GDNV\_\allowbreak{}BV} & \texttt{16} & \texttt{tools/\allowbreak{}gdn\_\allowbreak{}verify\_\allowbreak{}kernels.py} \\
\texttt{QWEN38\_\allowbreak{}GDNV\_\allowbreak{}CONV\_\allowbreak{}BLOCK} & \texttt{256} & \texttt{32} \textbullet{} \texttt{tools/\allowbreak{}gdn\_\allowbreak{}verify\_\allowbreak{}kernels.py} \\
\texttt{QWEN38\_\allowbreak{}GDNV\_\allowbreak{}CONV\_\allowbreak{}WARPS} & \texttt{4} & \texttt{1} \textbullet{} \texttt{tools/\allowbreak{}gdn\_\allowbreak{}verify\_\allowbreak{}kernels.py} \\
\texttt{QWEN38\_\allowbreak{}GDNV\_\allowbreak{}ONE\_\allowbreak{}WARP} & \texttt{0} & \texttt{1} \textbullet{} \texttt{tools/\allowbreak{}gdn\_\allowbreak{}verify\_\allowbreak{}kernels.py} \\
\texttt{QWEN38\_\allowbreak{}GDNV\_\allowbreak{}TREE\_\allowbreak{}PF} & \texttt{0} & \texttt{tools/\allowbreak{}gdn\_\allowbreak{}verify\_\allowbreak{}kernels.py} \\
\texttt{QWEN38\_\allowbreak{}GDNV\_\allowbreak{}WARPS} & \texttt{4} & \texttt{tools/\allowbreak{}gdn\_\allowbreak{}verify\_\allowbreak{}kernels.py} \\
\texttt{QWEN38\_\allowbreak{}GDNV\_\allowbreak{}WY} & \texttt{0} & \texttt{1} \textbullet{} \texttt{tools/\allowbreak{}gdn\_\allowbreak{}verify\_\allowbreak{}kernels.py} \\
\texttt{QWEN38\_\allowbreak{}GDNV\_\allowbreak{}WY\_\allowbreak{}CHAIN\_\allowbreak{}MAXT} & \texttt{32} & \texttt{16} \textbullet{} \texttt{tools/\allowbreak{}gdn\_\allowbreak{}verify\_\allowbreak{}kernels.py} \\
\texttt{QWEN38\_\allowbreak{}GDNV\_\allowbreak{}WY\_\allowbreak{}FUSED} & \texttt{0} & \texttt{1} \textbullet{} \texttt{tools/\allowbreak{}gdn\_\allowbreak{}verify\_\allowbreak{}kernels.py} \\
\texttt{QWEN38\_\allowbreak{}GDNV\_\allowbreak{}WY\_\allowbreak{}FUSED\_\allowbreak{}MAXT} & \texttt{32} & \texttt{16} \textbullet{} \texttt{tools/\allowbreak{}gdn\_\allowbreak{}verify\_\allowbreak{}kernels.py} \\
\texttt{QWEN38\_\allowbreak{}GDNV\_\allowbreak{}WY\_\allowbreak{}KC} & \texttt{0} & \texttt{tools/\allowbreak{}gdn\_\allowbreak{}verify\_\allowbreak{}kernels.py} \\
\texttt{QWEN38\_\allowbreak{}GDNV\_\allowbreak{}WY\_\allowbreak{}MAXT} & \texttt{32} & \texttt{16} \textbullet{} \texttt{tools/\allowbreak{}gdn\_\allowbreak{}verify\_\allowbreak{}kernels.py} \\
\texttt{QWEN38\_\allowbreak{}GDN\_\allowbreak{}AB} & \texttt{0} & \texttt{engine/\allowbreak{}model.py} \\
\texttt{QWEN38\_\allowbreak{}GDN\_\allowbreak{}CHUNK} & \texttt{64} & \texttt{engine/\allowbreak{}model.py} \\
\texttt{QWEN38\_\allowbreak{}GDN\_\allowbreak{}MM} & \texttt{fp32} & \texttt{engine/\allowbreak{}gdn.py} \\
\texttt{QWEN38\_\allowbreak{}GDN\_\allowbreak{}MM\_\allowbreak{}FROM} & \texttt{64} & \texttt{engine/\allowbreak{}gdn.py} \\
\texttt{QWEN38\_\allowbreak{}GDN\_\allowbreak{}PREFILL\_\allowbreak{}BV} & \texttt{32} & \texttt{tools/\allowbreak{}gdn\_\allowbreak{}prefill\_\allowbreak{}kernels.py} \\
\texttt{QWEN38\_\allowbreak{}GDN\_\allowbreak{}PREFILL\_\allowbreak{}NOREP} & \texttt{0} & \texttt{engine/\allowbreak{}model.py} \\
\texttt{QWEN38\_\allowbreak{}GDN\_\allowbreak{}PREFILL\_\allowbreak{}PREC} & \texttt{bf16x3} & \texttt{tools/\allowbreak{}gdn\_\allowbreak{}prefill\_\allowbreak{}kernels.py} \\
\texttt{QWEN38\_\allowbreak{}GDN\_\allowbreak{}PREFILL\_\allowbreak{}PREC\_\allowbreak{}A} & \texttt{tf32} & \texttt{tools/\allowbreak{}gdn\_\allowbreak{}prefill\_\allowbreak{}kernels.py} \\
\texttt{QWEN38\_\allowbreak{}GDN\_\allowbreak{}PREFILL\_\allowbreak{}PREC\_\allowbreak{}SCAN} & \texttt{none} & \texttt{tools/\allowbreak{}gdn\_\allowbreak{}prefill\_\allowbreak{}kernels.py} \\
\texttt{QWEN38\_\allowbreak{}GDN\_\allowbreak{}PREFILL\_\allowbreak{}SERIES} & \texttt{0} & \texttt{tools/\allowbreak{}gdn\_\allowbreak{}prefill\_\allowbreak{}kernels.py} \\
\texttt{QWEN38\_\allowbreak{}GDN\_\allowbreak{}PREFILL\_\allowbreak{}STAGES} & \texttt{1} & \texttt{tools/\allowbreak{}gdn\_\allowbreak{}prefill\_\allowbreak{}kernels.py} \\
\texttt{QWEN38\_\allowbreak{}GDN\_\allowbreak{}PREFILL\_\allowbreak{}WARPS\_\allowbreak{}INTRA} & \texttt{8} & \texttt{tools/\allowbreak{}gdn\_\allowbreak{}prefill\_\allowbreak{}kernels.py} \\
\texttt{QWEN38\_\allowbreak{}GDN\_\allowbreak{}PREFILL\_\allowbreak{}WARPS\_\allowbreak{}SCAN} & \texttt{8} & \texttt{tools/\allowbreak{}gdn\_\allowbreak{}prefill\_\allowbreak{}kernels.py} \\
\texttt{QWEN38\_\allowbreak{}GQA\_\allowbreak{}FROM} & \texttt{64} & \texttt{engine/\allowbreak{}model.py} \\
\texttt{QWEN38\_\allowbreak{}GRAPH\_\allowbreak{}MAX\_\allowbreak{}CTX} & \texttt{32768} & \texttt{engine/\allowbreak{}verify\_\allowbreak{}graph.py} \\
\texttt{QWEN38\_\allowbreak{}HEAD\_\allowbreak{}GEMM} & \texttt{none} & \texttt{tools/\allowbreak{}head\_\allowbreak{}gemv.py} \\
\texttt{QWEN38\_\allowbreak{}HOST\_\allowbreak{}ASYNC} & \texttt{0} & \texttt{engine/\allowbreak{}model.py} \\
\texttt{QWEN38\_\allowbreak{}HOST\_\allowbreak{}WALK} & \texttt{0} & \texttt{1} \textbullet{} \texttt{engine/\allowbreak{}drafters/\allowbreak{}dflash2.py} \\
\texttt{QWEN38\_\allowbreak{}KV\_\allowbreak{}FP8} & \texttt{0} & \texttt{engine/\allowbreak{}model.py} \\
\texttt{QWEN38\_\allowbreak{}LATCH\_\allowbreak{}PRICE} & \texttt{0} & \texttt{engine/\allowbreak{}lenrouter.py} \\
\texttt{QWEN38\_\allowbreak{}LAUNCH\_\allowbreak{}FIRST} & \texttt{0} & \texttt{server/\allowbreak{}app.py} \\
\texttt{QWEN38\_\allowbreak{}LEN\_\allowbreak{}MODE} & \texttt{none} & \texttt{wide} \textbullet{} \texttt{engine/\allowbreak{}lenrouter.py} \\
\texttt{QWEN38\_\allowbreak{}LEN\_\allowbreak{}SWITCH} & \texttt{0} & \texttt{1} \textbullet{} \texttt{engine/\allowbreak{}lenrouter.py} \\
\texttt{QWEN38\_\allowbreak{}LOWER\_\allowbreak{}RIGHT\_\allowbreak{}FROM} & \texttt{64} & \texttt{engine/\allowbreak{}model.py} \\
\texttt{QWEN38\_\allowbreak{}MEMORY\_\allowbreak{}PROXIES} & \texttt{none} & \texttt{server/\allowbreak{}app.py} \\
\texttt{QWEN38\_\allowbreak{}MEMORY\_\allowbreak{}WRITER\_\allowbreak{}TEXT} & \texttt{none} & \texttt{server/\allowbreak{}app.py} \\
\texttt{QWEN38\_\allowbreak{}MM\_\allowbreak{}END} & \texttt{on} & \texttt{server/\allowbreak{}app.py} \\
\texttt{QWEN38\_\allowbreak{}MODEL} & \texttt{none} & \texttt{engine/\allowbreak{}config.py} \\
\texttt{QWEN38\_\allowbreak{}MODEL\_\allowbreak{}NAME} & \texttt{none} & \texttt{server/\allowbreak{}conductor/\allowbreak{}schema.py} \\
\texttt{QWEN38\_\allowbreak{}NARRATOR\_\allowbreak{}DEVICE} & \texttt{auto} & \texttt{server/\allowbreak{}conductor/\allowbreak{}narrator.py} \\
\texttt{QWEN38\_\allowbreak{}NARRATOR\_\allowbreak{}DIR} & \texttt{none} & \texttt{server/\allowbreak{}conductor/\allowbreak{}narrator.py} \\
\texttt{QWEN38\_\allowbreak{}NARRATOR\_\allowbreak{}MODEL} & \texttt{none} & \texttt{server/\allowbreak{}conductor/\allowbreak{}narrator.py} \\
\texttt{QWEN38\_\allowbreak{}NOTE\_\allowbreak{}ANCHOR} & \texttt{on} & \texttt{server/\allowbreak{}conductor/\allowbreak{}router.py} \\
\texttt{QWEN38\_\allowbreak{}NOTE\_\allowbreak{}DEDUP} & \texttt{parts} & \texttt{server/\allowbreak{}conductor/\allowbreak{}router.py} \\
\texttt{QWEN38\_\allowbreak{}NVFP4} & \texttt{none} & \texttt{engine/\allowbreak{}loader.py} \\
\texttt{QWEN38\_\allowbreak{}NVFP4\_\allowbreak{}BLOCK\_\allowbreak{}TILES} & \texttt{0} & \texttt{tools/\allowbreak{}nvfp4\_\allowbreak{}linear.py} \\
\texttt{QWEN38\_\allowbreak{}NVFP4\_\allowbreak{}DEQUANT\_\allowbreak{}FROM} & \texttt{512} & \texttt{tools/\allowbreak{}nvfp4\_\allowbreak{}linear.py} \\
\texttt{QWEN38\_\allowbreak{}NVFP4\_\allowbreak{}PREFILL\_\allowbreak{}V2} & \texttt{0} & \texttt{1} \textbullet{} \texttt{tools/\allowbreak{}nvfp4\_\allowbreak{}linear.py} \\
\texttt{QWEN38\_\allowbreak{}NVFP4\_\allowbreak{}PREFILL\_\allowbreak{}V2\_\allowbreak{}UNTIL} & \texttt{1024} & \texttt{1025} \textbullet{} \texttt{tools/\allowbreak{}nvfp4\_\allowbreak{}linear.py} \\
\texttt{QWEN38\_\allowbreak{}NVFP4\_\allowbreak{}SKINNY} & \texttt{0} & \texttt{1} \textbullet{} \texttt{tools/\allowbreak{}nvfp4\_\allowbreak{}skinny.py} \\
\texttt{QWEN38\_\allowbreak{}NVFP4\_\allowbreak{}SPLITK} & \texttt{0} & \texttt{tools/\allowbreak{}nvfp4\_\allowbreak{}linear.py} \\
\texttt{QWEN38\_\allowbreak{}NVFP4\_\allowbreak{}V2} & \texttt{1} & \texttt{tools/\allowbreak{}nvfp4\_\allowbreak{}linear\_\allowbreak{}v2.py} \\
\texttt{QWEN38\_\allowbreak{}NVFP4\_\allowbreak{}V2\_\allowbreak{}BN} & \texttt{0} & \texttt{tools/\allowbreak{}nvfp4\_\allowbreak{}linear\_\allowbreak{}v2.py} \\
\texttt{QWEN38\_\allowbreak{}NVFP4\_\allowbreak{}V2\_\allowbreak{}DOTS} & \texttt{1} & \texttt{tools/\allowbreak{}nvfp4\_\allowbreak{}linear\_\allowbreak{}v2.py} \\
\texttt{QWEN38\_\allowbreak{}NVFP4\_\allowbreak{}V2\_\allowbreak{}MAX} & \texttt{32} & \texttt{tools/\allowbreak{}nvfp4\_\allowbreak{}linear\_\allowbreak{}v2.py} \\
\texttt{QWEN38\_\allowbreak{}NVFP4\_\allowbreak{}V2\_\allowbreak{}MIN} & \texttt{1} & \texttt{tools/\allowbreak{}nvfp4\_\allowbreak{}linear\_\allowbreak{}v2.py} \\
\texttt{QWEN38\_\allowbreak{}NVFP4\_\allowbreak{}V2\_\allowbreak{}PREFETCH} & \texttt{0} & \texttt{tools/\allowbreak{}nvfp4\_\allowbreak{}linear\_\allowbreak{}v2.py} \\
\texttt{QWEN38\_\allowbreak{}NVFP4\_\allowbreak{}V2\_\allowbreak{}W} & \texttt{0} & \texttt{tools/\allowbreak{}nvfp4\_\allowbreak{}linear\_\allowbreak{}v2.py} \\
\texttt{QWEN38\_\allowbreak{}PREFILL\_\allowbreak{}CAUSAL} & \texttt{1} & \texttt{engine/\allowbreak{}model.py} \\
\texttt{QWEN38\_\allowbreak{}PREFILL\_\allowbreak{}WINDOW} & \texttt{8192} & \texttt{engine/\allowbreak{}cache.py} \\
\texttt{QWEN38\_\allowbreak{}PREFILL\_\allowbreak{}WINDOW\_\allowbreak{}FREE} & \texttt{32768} & \texttt{engine/\allowbreak{}cache.py} \\
\texttt{QWEN38\_\allowbreak{}PROBE\_\allowbreak{}DAY\_\allowbreak{}MB} & \texttt{100} & \texttt{server/\allowbreak{}probe.py} \\
\texttt{QWEN38\_\allowbreak{}PROBE\_\allowbreak{}DIR} & \texttt{\textasciitilde{}/\allowbreak{}.qwen38-\allowbreak{}spark-\allowbreak{}engine/\allowbreak{}probe} & \texttt{(a path)} \textbullet{} \texttt{server/\allowbreak{}probe.py} \\
\texttt{QWEN38\_\allowbreak{}PROBE\_\allowbreak{}EVERY} & \texttt{1024} & \texttt{server/\allowbreak{}probe.py} \\
\texttt{QWEN38\_\allowbreak{}PROBE\_\allowbreak{}SNAPSHOTS} & \texttt{off} & \texttt{off} \textbullet{} \texttt{server/\allowbreak{}probe.py} \\
\texttt{QWEN38\_\allowbreak{}Q35M\_\allowbreak{}ANCHOR\_\allowbreak{}BACKOFF} & \texttt{0} & \texttt{engine/\allowbreak{}models/\allowbreak{}qwen3\_\allowbreak{}5\_\allowbreak{}moe/\allowbreak{}serving.py} \\
\texttt{QWEN38\_\allowbreak{}Q35M\_\allowbreak{}DFLASH2\_\allowbreak{}HOT} & \texttt{0} & \texttt{engine/\allowbreak{}models/\allowbreak{}qwen3\_\allowbreak{}5\_\allowbreak{}moe/\allowbreak{}dflash2.py} \\
\texttt{QWEN38\_\allowbreak{}Q35M\_\allowbreak{}DFLASH2\_\allowbreak{}NVFP4} & \texttt{0} & \texttt{engine/\allowbreak{}models/\allowbreak{}qwen3\_\allowbreak{}5\_\allowbreak{}moe/\allowbreak{}dflash2.py} \\
\texttt{QWEN38\_\allowbreak{}Q35M\_\allowbreak{}LOOKUP} & \texttt{1} & \texttt{engine/\allowbreak{}models/\allowbreak{}qwen3\_\allowbreak{}5\_\allowbreak{}moe/\allowbreak{}serving.py} \\
\texttt{QWEN38\_\allowbreak{}Q35M\_\allowbreak{}SIBLINGS} & \texttt{0} & \texttt{engine/\allowbreak{}models/\allowbreak{}qwen3\_\allowbreak{}5\_\allowbreak{}moe/\allowbreak{}serving.py} \\
\texttt{QWEN38\_\allowbreak{}RANKK} & \texttt{1} & \texttt{engine/\allowbreak{}model.py} \\
\texttt{QWEN38\_\allowbreak{}REGENERATE} & \texttt{sample} & \texttt{server/\allowbreak{}app.py} \\
\texttt{QWEN38\_\allowbreak{}REGEN\_\allowbreak{}TEMPERATURE} & \texttt{0.7} & \texttt{server/\allowbreak{}app.py} \\
\texttt{QWEN38\_\allowbreak{}REGEN\_\allowbreak{}TOP\_\allowbreak{}K} & \texttt{20} & \texttt{server/\allowbreak{}app.py} \\
\texttt{QWEN38\_\allowbreak{}REGEN\_\allowbreak{}TOP\_\allowbreak{}P} & \texttt{0.95} & \texttt{server/\allowbreak{}app.py} \\
\texttt{QWEN38\_\allowbreak{}REREAD\_\allowbreak{}SAY\_\allowbreak{}TOKENS} & \texttt{8192} & \texttt{server/\allowbreak{}app.py} \\
\texttt{QWEN38\_\allowbreak{}SAMPLED\_\allowbreak{}CUT} & \texttt{0} & \texttt{engine/\allowbreak{}router.py} \\
\texttt{QWEN38\_\allowbreak{}SAMPLE\_\allowbreak{}BATCH} & \texttt{1} & \texttt{engine/\allowbreak{}sample.py} \\
\texttt{QWEN38\_\allowbreak{}SAMPLE\_\allowbreak{}NUCLEUS} & \texttt{1} & \texttt{engine/\allowbreak{}sample.py} \\
\texttt{QWEN38\_\allowbreak{}SAMPLE\_\allowbreak{}TOPP\_\allowbreak{}K} & \texttt{1024} & \texttt{engine/\allowbreak{}sample.py} \\
\texttt{QWEN38\_\allowbreak{}SCALE\_\allowbreak{}ON} & \texttt{weight} & \texttt{tools/\allowbreak{}fp8\_\allowbreak{}linear.py} \\
\texttt{QWEN38\_\allowbreak{}SESSION\_\allowbreak{}GENERATIONS} & \texttt{on} & \texttt{server/\allowbreak{}conductor/\allowbreak{}router.py} \\
\texttt{QWEN38\_\allowbreak{}SHADOW\_\allowbreak{}DIR} & \texttt{none} & \texttt{server/\allowbreak{}conductor/\allowbreak{}shadow.py} \\
\texttt{QWEN38\_\allowbreak{}SHADOW\_\allowbreak{}DISK} & \texttt{none} & \texttt{auto} \textbullet{} \texttt{server/\allowbreak{}conductor/\allowbreak{}shadow.py} \\
\texttt{QWEN38\_\allowbreak{}SHADOW\_\allowbreak{}DISK\_\allowbreak{}GB} & \texttt{20} & \texttt{20} \textbullet{} \texttt{server/\allowbreak{}conductor/\allowbreak{}shadow.py} \\
\texttt{QWEN38\_\allowbreak{}SHADOW\_\allowbreak{}GB} & \texttt{4} & \texttt{4} \textbullet{} \texttt{server/\allowbreak{}conductor/\allowbreak{}shadow.py} \\
\texttt{QWEN38\_\allowbreak{}SHADOW\_\allowbreak{}GRID} & \texttt{512} & \texttt{512} \textbullet{} \texttt{server/\allowbreak{}conductor/\allowbreak{}shadow.py} \\
\texttt{QWEN38\_\allowbreak{}SHADOW\_\allowbreak{}MODEL} & \texttt{none} & \texttt{server/\allowbreak{}conductor/\allowbreak{}shadow.py} \\
\texttt{QWEN38\_\allowbreak{}SHADOW\_\allowbreak{}POOL\_\allowbreak{}SLACK\_\allowbreak{}GIB} & \texttt{2} & \texttt{server/\allowbreak{}conductor/\allowbreak{}shadow.py} \\
\texttt{QWEN38\_\allowbreak{}SIDE\_\allowbreak{}CARRY} & \texttt{on} & \texttt{server/\allowbreak{}app.py} \\
\texttt{QWEN38\_\allowbreak{}SIDE\_\allowbreak{}CONTINUE} & \texttt{view} & \texttt{server/\allowbreak{}app.py} \\
\texttt{QWEN38\_\allowbreak{}SIDE\_\allowbreak{}CONV} & \texttt{on} & \texttt{server/\allowbreak{}app.py} \\
\texttt{QWEN38\_\allowbreak{}SIDE\_\allowbreak{}DISK} & \texttt{read} & \texttt{server/\allowbreak{}app.py} \\
\texttt{QWEN38\_\allowbreak{}SIDE\_\allowbreak{}GUESTS} & \texttt{on} & \texttt{server/\allowbreak{}app.py} \\
\texttt{QWEN38\_\allowbreak{}SIDE\_\allowbreak{}HOME\_\allowbreak{}WAIT\_\allowbreak{}S} & \texttt{10} & \texttt{server/\allowbreak{}app.py} \\
\texttt{QWEN38\_\allowbreak{}SIDE\_\allowbreak{}IMAGES} & \texttt{on} & \texttt{server/\allowbreak{}app.py} \\
\texttt{QWEN38\_\allowbreak{}SIDE\_\allowbreak{}MIN\_\allowbreak{}ROOM} & \texttt{16384} & \texttt{server/\allowbreak{}app.py} \\
\texttt{QWEN38\_\allowbreak{}SIDE\_\allowbreak{}NEW\_\allowbreak{}FIRST} & \texttt{on} & \texttt{server/\allowbreak{}app.py} \\
\texttt{QWEN38\_\allowbreak{}SIDE\_\allowbreak{}NEW\_\allowbreak{}WAIT\_\allowbreak{}S} & \texttt{20} & \texttt{server/\allowbreak{}app.py} \\
\texttt{QWEN38\_\allowbreak{}SIDE\_\allowbreak{}RESIDENT\_\allowbreak{}GB} & \texttt{0.5} & \texttt{server/\allowbreak{}app.py} \\
\texttt{QWEN38\_\allowbreak{}SIDE\_\allowbreak{}SLOT0\_\allowbreak{}WAIT\_\allowbreak{}S} & \texttt{3} & \texttt{server/\allowbreak{}app.py} \\
\texttt{QWEN38\_\allowbreak{}SIDE\_\allowbreak{}STASH\_\allowbreak{}GB} & \texttt{1.0} & \texttt{server/\allowbreak{}app.py} \\
\texttt{QWEN38\_\allowbreak{}SKINNY\_\allowbreak{}LDW} & \texttt{0} & \texttt{tools/\allowbreak{}nvfp4\_\allowbreak{}skinny.py} \\
\texttt{QWEN38\_\allowbreak{}SKINNY\_\allowbreak{}PDL} & \texttt{0} & \texttt{tools/\allowbreak{}nvfp4\_\allowbreak{}skinny.py} \\
\texttt{QWEN38\_\allowbreak{}SKINNY\_\allowbreak{}SRUN} & \texttt{0} & \texttt{1} \textbullet{} \texttt{tools/\allowbreak{}nvfp4\_\allowbreak{}skinny.py} \\
\texttt{QWEN38\_\allowbreak{}SKINNY\_\allowbreak{}SSTUB} & \texttt{0} & \texttt{tools/\allowbreak{}nvfp4\_\allowbreak{}skinny.py} \\
\texttt{QWEN38\_\allowbreak{}SKINNY\_\allowbreak{}TILES} & \texttt{none} & \texttt{(a path)} \textbullet{} \texttt{tools/\allowbreak{}nvfp4\_\allowbreak{}skinny.py} \\
\texttt{QWEN38\_\allowbreak{}SKINNY\_\allowbreak{}TILES\_\allowbreak{}B} & \texttt{none} & \texttt{tools/\allowbreak{}nvfp4\_\allowbreak{}skinny.py} \\
\texttt{QWEN38\_\allowbreak{}SKINNY\_\allowbreak{}TILES\_\allowbreak{}C} & \texttt{none} & \texttt{tools/\allowbreak{}nvfp4\_\allowbreak{}skinny.py} \\
\texttt{QWEN38\_\allowbreak{}SKINNY\_\allowbreak{}TILES\_\allowbreak{}WIDE} & \texttt{none} & \texttt{(a path)} \textbullet{} \texttt{tools/\allowbreak{}nvfp4\_\allowbreak{}skinny.py} \\
\texttt{QWEN38\_\allowbreak{}SKINNY\_\allowbreak{}TILES\_\allowbreak{}WIDE\_\allowbreak{}B} & \texttt{none} & \texttt{tools/\allowbreak{}nvfp4\_\allowbreak{}skinny.py} \\
\texttt{QWEN38\_\allowbreak{}SKINNY\_\allowbreak{}XSTUB} & \texttt{0} & \texttt{tools/\allowbreak{}nvfp4\_\allowbreak{}skinny.py} \\
\texttt{QWEN38\_\allowbreak{}SLOT0\_\allowbreak{}KEEP\_\allowbreak{}S} & \texttt{1800} & \texttt{server/\allowbreak{}app.py} \\
\texttt{QWEN38\_\allowbreak{}SPILL\_\allowbreak{}PARK} & \texttt{on} & \texttt{engine/\allowbreak{}cache.py} \\
\texttt{QWEN38\_\allowbreak{}STAIR\_\allowbreak{}MAX\_\allowbreak{}NODES} & \texttt{none} & \texttt{engine/\allowbreak{}lenrouter.py} \\
\texttt{QWEN38\_\allowbreak{}STAIR\_\allowbreak{}RHO} & \texttt{none} & \texttt{(a path)} \textbullet{} \texttt{engine/\allowbreak{}lenrouter.py} \\
\texttt{QWEN38\_\allowbreak{}STAIR\_\allowbreak{}SKIP} & \texttt{0} & \texttt{1} \textbullet{} \texttt{engine/\allowbreak{}lenrouter.py} \\
\texttt{QWEN38\_\allowbreak{}STAIR\_\allowbreak{}TABLES} & \texttt{none} & \texttt{(a path)} \textbullet{} \texttt{engine/\allowbreak{}lenrouter.py} \\
\texttt{QWEN38\_\allowbreak{}STATUS\_\allowbreak{}CLIENTS} & \texttt{open-\allowbreak{}webui} & \texttt{server/\allowbreak{}conductor/\allowbreak{}states.py} \\
\texttt{QWEN38\_\allowbreak{}SUFFIX\_\allowbreak{}STORE} & \texttt{\textasciitilde{}/\allowbreak{}.qwen38-\allowbreak{}spark-\allowbreak{}engine/\allowbreak{}suffix} & \texttt{server/\allowbreak{}app.py} \\
\texttt{QWEN38\_\allowbreak{}TANDEM\_\allowbreak{}SCRATCH} & \texttt{none} & \texttt{server/\allowbreak{}conductor/\allowbreak{}packages.py} \\
\texttt{QWEN38\_\allowbreak{}TANDEM\_\allowbreak{}URL} & \texttt{none} & \texttt{server/\allowbreak{}conductor/\allowbreak{}packages.py} \\
\texttt{QWEN38\_\allowbreak{}TOOL\_\allowbreak{}FIT} & \texttt{on} & \texttt{server/\allowbreak{}conductor/\allowbreak{}router.py} \\
\texttt{QWEN38\_\allowbreak{}TREE\_\allowbreak{}ALIAS\_\allowbreak{}STATE} & \texttt{0} & \texttt{0} \textbullet{} \texttt{engine/\allowbreak{}model.py} \\
\texttt{QWEN38\_\allowbreak{}TREE\_\allowbreak{}CHAIN\_\allowbreak{}DELEGATE} & \texttt{1} & \texttt{engine/\allowbreak{}model.py} \\
\texttt{QWEN38\_\allowbreak{}TREE\_\allowbreak{}HOST\_\allowbreak{}DEPTH} & \texttt{0} & \texttt{1} \textbullet{} \texttt{engine/\allowbreak{}model.py} \\
\texttt{QWEN38\_\allowbreak{}TREE\_\allowbreak{}MS} & \texttt{none} & \texttt{8:\allowbreak{}77.9,\allowbreak{}16:\allowbreak{}78.9,\allowbreak{}24:\allowbreak{}87.0,\allowbreak{}32:\allowbreak{}97.0} \textbullet{} \texttt{engine/\allowbreak{}router.py} \\
\texttt{QWEN38\_\allowbreak{}TREE\_\allowbreak{}NODES} & \texttt{0} & \texttt{24} \textbullet{} \texttt{engine/\allowbreak{}router.py} \\
\texttt{QWEN38\_\allowbreak{}TREE\_\allowbreak{}NODES\_\allowbreak{}NARROW} & \texttt{0} & \texttt{16} \textbullet{} \texttt{engine/\allowbreak{}router.py} \\
\texttt{QWEN38\_\allowbreak{}TREE\_\allowbreak{}WIDE\_\allowbreak{}AFTER} & \texttt{0} & \texttt{32} \textbullet{} \texttt{engine/\allowbreak{}lenrouter.py} \\
\texttt{QWEN38\_\allowbreak{}TWO\_\allowbreak{}STREAM} & \texttt{0} & \texttt{engine/\allowbreak{}model.py} \\
\texttt{QWEN38\_\allowbreak{}UT\_\allowbreak{}INVERSE} & \texttt{1} & \texttt{engine/\allowbreak{}gdn.py} \\
\texttt{QWEN38\_\allowbreak{}VERIFY\_\allowbreak{}GRAPH} & \texttt{0} & \texttt{1} \textbullet{} \texttt{engine/\allowbreak{}model.py} \\
\texttt{QWEN38\_\allowbreak{}VERIFY\_\allowbreak{}ROWS} & \texttt{16} & \texttt{32} \textbullet{} \texttt{engine/\allowbreak{}model.py} \\
\texttt{QWEN38\_\allowbreak{}VERIFY\_\allowbreak{}TILES} & \texttt{1} & \texttt{tools/\allowbreak{}nvfp4\_\allowbreak{}verify\_\allowbreak{}tiles.py} \\
\texttt{QWEN38\_\allowbreak{}WARM\_\allowbreak{}PREFILL\_\allowbreak{}TOKENS} & \texttt{2600} & \texttt{server/\allowbreak{}app.py} \\
\texttt{QWEN38\_\allowbreak{}WY\_\allowbreak{}BLOCK} & \texttt{8} & \texttt{tools/\allowbreak{}gdn\_\allowbreak{}wy\_\allowbreak{}kernels.py} \\
\texttt{QWEN38\_\allowbreak{}WY\_\allowbreak{}BV} & \texttt{32} & \texttt{tools/\allowbreak{}gdn\_\allowbreak{}wy\_\allowbreak{}kernels.py} \\
\texttt{QWEN38\_\allowbreak{}WY\_\allowbreak{}DBG} & \texttt{0} & \texttt{tools/\allowbreak{}gdn\_\allowbreak{}wy\_\allowbreak{}kernels.py} \\
\texttt{QWEN38\_\allowbreak{}WY\_\allowbreak{}PREC} & \texttt{ieee} & \texttt{tools/\allowbreak{}gdn\_\allowbreak{}wy\_\allowbreak{}kernels.py} \\
\texttt{QWEN38\_\allowbreak{}WY\_\allowbreak{}WARPS} & \texttt{4} & \texttt{tools/\allowbreak{}gdn\_\allowbreak{}wy\_\allowbreak{}kernels.py} \\
\texttt{QWEN38\_\allowbreak{}WY\_\allowbreak{}WARPS\_\allowbreak{}PREP} & \texttt{4} & \texttt{tools/\allowbreak{}gdn\_\allowbreak{}wy\_\allowbreak{}kernels.py} \\
\end{reftable}


\section{Server variables}\index{QSE variables}

% Generated by manual/tools/gen_reference.py from the code. Do not edit by hand.

\begin{reftabletwo}{QSE\_* variables: the server's settings from the environment or serve.env}{tab:ref-qse}{\hd{Variable} & \hd{Read in}}
\texttt{QSE\_\allowbreak{}ADMIN\_\allowbreak{}TOKEN} & \texttt{ops/\allowbreak{}chatd.service,\allowbreak{} ops/\allowbreak{}chatd.sh,\allowbreak{} ops/\allowbreak{}make-\allowbreak{}secrets.sh and 9 more} \\
\texttt{QSE\_\allowbreak{}ANSWER\_\allowbreak{}RESERVE} & \texttt{server/\allowbreak{}app.py} \\
\texttt{QSE\_\allowbreak{}ASSETS} & \texttt{ops/\allowbreak{}profiles/\allowbreak{}\_\allowbreak{}flashnext.env,\allowbreak{} ops/\allowbreak{}profiles/\allowbreak{}\_\allowbreak{}lfm25.env,\allowbreak{} ops/\allowbreak{}profiles/\allowbreak{}\_\allowbreak{}nemotron.env and 15 more} \\
\texttt{QSE\_\allowbreak{}BACKENDS} & \texttt{ops/\allowbreak{}serve.env,\allowbreak{} ops/\allowbreak{}start.sh,\allowbreak{} server/\allowbreak{}app.py and 1 more} \\
\texttt{QSE\_\allowbreak{}BARRIER\_\allowbreak{}IDLE\_\allowbreak{}OWNER} & \texttt{ops/\allowbreak{}serve.env,\allowbreak{} server/\allowbreak{}slots.py} \\
\texttt{QSE\_\allowbreak{}CHATD} & \texttt{install.sh,\allowbreak{} ops/\allowbreak{}chatd.service,\allowbreak{} ops/\allowbreak{}start.sh and 3 more} \\
\texttt{QSE\_\allowbreak{}CHATD\_\allowbreak{}ENV} & \texttt{ops/\allowbreak{}chatd.sh} \\
\texttt{QSE\_\allowbreak{}CHAT\_\allowbreak{}DATA} & \texttt{ops/\allowbreak{}chatd.sh,\allowbreak{} ops/\allowbreak{}sandbox/\allowbreak{}uninstall.sh,\allowbreak{} server/\allowbreak{}chat\_\allowbreak{}forward.py and 1 more} \\
\texttt{QSE\_\allowbreak{}CHAT\_\allowbreak{}ENGINE} & \texttt{ops/\allowbreak{}chatd.sh} \\
\texttt{QSE\_\allowbreak{}CHAT\_\allowbreak{}FORWARD\_\allowbreak{}KEY} & \texttt{ops/\allowbreak{}chatd.sh,\allowbreak{} ops/\allowbreak{}make-\allowbreak{}secrets.sh,\allowbreak{} server/\allowbreak{}chat\_\allowbreak{}forward.py} \\
\texttt{QSE\_\allowbreak{}CHAT\_\allowbreak{}URL} & \texttt{ops/\allowbreak{}chatd.sh,\allowbreak{} server/\allowbreak{}chat\_\allowbreak{}forward.py} \\
\texttt{QSE\_\allowbreak{}COMMUNITY\_\allowbreak{}DIR} & \texttt{server/\allowbreak{}setup/\allowbreak{}community.py} \\
\texttt{QSE\_\allowbreak{}COMMUNITY\_\allowbreak{}ENV\_\allowbreak{}DIR} & \texttt{ops/\allowbreak{}serve.env,\allowbreak{} ops/\allowbreak{}use-\allowbreak{}profile.sh} \\
\texttt{QSE\_\allowbreak{}COMMUNITY\_\allowbreak{}URL} & \texttt{ops/\allowbreak{}serve.env,\allowbreak{} server/\allowbreak{}setup/\allowbreak{}community.py} \\
\texttt{QSE\_\allowbreak{}COMPACTION} & \texttt{server/\allowbreak{}app.py} \\
\texttt{QSE\_\allowbreak{}COMPACTION\_\allowbreak{}IDLE\_\allowbreak{}S} & \texttt{server/\allowbreak{}app.py} \\
\texttt{QSE\_\allowbreak{}COMPACTION\_\allowbreak{}MIN\_\allowbreak{}TOKENS} & \texttt{server/\allowbreak{}app.py} \\
\texttt{QSE\_\allowbreak{}COMPACTION\_\allowbreak{}PREPARE} & \texttt{server/\allowbreak{}app.py} \\
\texttt{QSE\_\allowbreak{}CONDUCTOR\_\allowbreak{}CONFIG} & \texttt{ops/\allowbreak{}serve.env,\allowbreak{} ops/\allowbreak{}start.sh,\allowbreak{} ops/\allowbreak{}use-\allowbreak{}profile.sh and 3 more} \\
\texttt{QSE\_\allowbreak{}CONDUCTOR\_\allowbreak{}ROOT} & \texttt{ops/\allowbreak{}roles.sh,\allowbreak{} server/\allowbreak{}conductor/\allowbreak{}residency.py} \\
\texttt{QSE\_\allowbreak{}DASHBOARD\_\allowbreak{}LOGIN} & \texttt{ops/\allowbreak{}serve.env,\allowbreak{} ops/\allowbreak{}start.sh,\allowbreak{} proxy/\allowbreak{}\_\allowbreak{}\_\allowbreak{}main\_\allowbreak{}\_\allowbreak{}.py and 2 more} \\
\texttt{QSE\_\allowbreak{}DISK\_\allowbreak{}ADOPT\_\allowbreak{}FROM} & \texttt{ops/\allowbreak{}serve.env,\allowbreak{} server/\allowbreak{}diskcache.py} \\
\texttt{QSE\_\allowbreak{}DISK\_\allowbreak{}CACHE} & \texttt{ops/\allowbreak{}profiles/\allowbreak{}\_\allowbreak{}exl3.env,\allowbreak{} ops/\allowbreak{}profiles/\allowbreak{}\_\allowbreak{}flashnext.env,\allowbreak{} ops/\allowbreak{}profiles/\allowbreak{}\_\allowbreak{}lfm25.env and 8 more} \\
\texttt{QSE\_\allowbreak{}ENGINE\_\allowbreak{}FLAGS} & \texttt{ops/\allowbreak{}profiles/\allowbreak{}\_\allowbreak{}flashnext.env,\allowbreak{} ops/\allowbreak{}profiles/\allowbreak{}\_\allowbreak{}lfm25.env,\allowbreak{} ops/\allowbreak{}profiles/\allowbreak{}\_\allowbreak{}nemotron.env and 4 more} \\
\texttt{QSE\_\allowbreak{}ENGINE\_\allowbreak{}PORT} & \texttt{ops/\allowbreak{}engine\_\allowbreak{}status.py} \\
\texttt{QSE\_\allowbreak{}EXL3\_\allowbreak{}BPW} & \texttt{ops/\allowbreak{}profiles/\allowbreak{}\_\allowbreak{}exl3.env,\allowbreak{} ops/\allowbreak{}profiles/\allowbreak{}exl3-\allowbreak{}2.5.env,\allowbreak{} ops/\allowbreak{}profiles/\allowbreak{}exl3-\allowbreak{}3.0.env and 4 more} \\
\texttt{QSE\_\allowbreak{}FAKE\_\allowbreak{}CARD\_\allowbreak{}DIR} & \texttt{server/\allowbreak{}fake\_\allowbreak{}engine.py} \\
\texttt{QSE\_\allowbreak{}FAKE\_\allowbreak{}PROMPTS} & \texttt{server/\allowbreak{}fake\_\allowbreak{}engine.py} \\
\texttt{QSE\_\allowbreak{}FETCH\_\allowbreak{}ASSETS} & \texttt{ops/\allowbreak{}use-\allowbreak{}profile.sh,\allowbreak{} server/\allowbreak{}setup/\allowbreak{}apply.py} \\
\texttt{QSE\_\allowbreak{}FN\_\allowbreak{}BPW} & \texttt{ops/\allowbreak{}profiles/\allowbreak{}\_\allowbreak{}flashnext.env,\allowbreak{} ops/\allowbreak{}profiles/\allowbreak{}flashnext-\allowbreak{}2.05.env,\allowbreak{} ops/\allowbreak{}profiles/\allowbreak{}flashnext-\allowbreak{}3.05.env and 3 more} \\
\texttt{QSE\_\allowbreak{}FN\_\allowbreak{}TABLE} & \texttt{ops/\allowbreak{}profiles/\allowbreak{}\_\allowbreak{}flashnext.env,\allowbreak{} ops/\allowbreak{}profiles/\allowbreak{}flashnext-\allowbreak{}2.05.env,\allowbreak{} ops/\allowbreak{}profiles/\allowbreak{}flashnext-\allowbreak{}3.05.env and 3 more} \\
\texttt{QSE\_\allowbreak{}FN\_\allowbreak{}VISION\_\allowbreak{}FILE} & \texttt{ops/\allowbreak{}profiles/\allowbreak{}\_\allowbreak{}flashnext.env,\allowbreak{} ops/\allowbreak{}profiles/\allowbreak{}flashnext-\allowbreak{}4.05.env} \\
\texttt{QSE\_\allowbreak{}IMAGE\_\allowbreak{}BACKEND} & \texttt{server/\allowbreak{}app.py} \\
\texttt{QSE\_\allowbreak{}KV\_\allowbreak{}HARD\_\allowbreak{}FLOOR\_\allowbreak{}GB} & \texttt{ops/\allowbreak{}profiles/\allowbreak{}\_\allowbreak{}flashnext.env,\allowbreak{} ops/\allowbreak{}serve.env,\allowbreak{} server/\allowbreak{}app.py} \\
\texttt{QSE\_\allowbreak{}KV\_\allowbreak{}MARGIN\_\allowbreak{}ROWS} & \texttt{server/\allowbreak{}app.py} \\
\texttt{QSE\_\allowbreak{}KV\_\allowbreak{}POOL} & \texttt{engine/\allowbreak{}kvpool.py,\allowbreak{} ops/\allowbreak{}profiles/\allowbreak{}\_\allowbreak{}flashnext.env,\allowbreak{} ops/\allowbreak{}profiles/\allowbreak{}\_\allowbreak{}lfm25.env and 5 more} \\
\texttt{QSE\_\allowbreak{}KV\_\allowbreak{}POOL\_\allowbreak{}MIN\_\allowbreak{}ROWS} & \texttt{engine/\allowbreak{}kvpool.py} \\
\texttt{QSE\_\allowbreak{}LFM\_\allowbreak{}VARIANT} & \texttt{ops/\allowbreak{}profiles/\allowbreak{}\_\allowbreak{}lfm25.env,\allowbreak{} ops/\allowbreak{}profiles/\allowbreak{}lfm25-\allowbreak{}bf16.env,\allowbreak{} ops/\allowbreak{}profiles/\allowbreak{}lfm25-\allowbreak{}exl3-\allowbreak{}2.5.env and 5 more} \\
\texttt{QSE\_\allowbreak{}LIVE\_\allowbreak{}DECODE\_\allowbreak{}WARMUP\_\allowbreak{}TOKENS} & \texttt{ops/\allowbreak{}serve.env,\allowbreak{} server/\allowbreak{}live.py} \\
\texttt{QSE\_\allowbreak{}LIVE\_\allowbreak{}HZ} & \texttt{ops/\allowbreak{}serve.env,\allowbreak{} ops/\allowbreak{}start.sh,\allowbreak{} server/\allowbreak{}app.py and 1 more} \\
\texttt{QSE\_\allowbreak{}LIVE\_\allowbreak{}MOSTLY\_\allowbreak{}CACHED\_\allowbreak{}PCT} & \texttt{ops/\allowbreak{}serve.env,\allowbreak{} server/\allowbreak{}live.py} \\
\texttt{QSE\_\allowbreak{}LIVE\_\allowbreak{}MOSTLY\_\allowbreak{}CACHED\_\allowbreak{}TOKENS} & \texttt{ops/\allowbreak{}serve.env,\allowbreak{} server/\allowbreak{}live.py} \\
\texttt{QSE\_\allowbreak{}LIVE\_\allowbreak{}NOW\_\allowbreak{}WINDOW\_\allowbreak{}S} & \texttt{ops/\allowbreak{}serve.env,\allowbreak{} server/\allowbreak{}live.py} \\
\texttt{QSE\_\allowbreak{}LIVE\_\allowbreak{}PREFILL\_\allowbreak{}WARMUP\_\allowbreak{}CHUNKS} & \texttt{ops/\allowbreak{}serve.env,\allowbreak{} server/\allowbreak{}live.py} \\
\texttt{QSE\_\allowbreak{}MEM\_\allowbreak{}NEED\_\allowbreak{}GIB} & \texttt{ops/\allowbreak{}profiles/\allowbreak{}flashnext-\allowbreak{}2.05.env,\allowbreak{} ops/\allowbreak{}profiles/\allowbreak{}flashnext-\allowbreak{}3.05.env,\allowbreak{} ops/\allowbreak{}profiles/\allowbreak{}flashnext-\allowbreak{}4.05.env and 24 more} \\
\texttt{QSE\_\allowbreak{}METRICS\_\allowbreak{}TOKEN} & \texttt{ops/\allowbreak{}chatd.service,\allowbreak{} ops/\allowbreak{}chatd.sh,\allowbreak{} ops/\allowbreak{}make-\allowbreak{}secrets.sh and 3 more} \\
\texttt{QSE\_\allowbreak{}NH\_\allowbreak{}DENSE\_\allowbreak{}ROWS} & \texttt{engine/\allowbreak{}models/\allowbreak{}nemotron\_\allowbreak{}h/\allowbreak{}engine.py} \\
\texttt{QSE\_\allowbreak{}NH\_\allowbreak{}DIR} & \texttt{ops/\allowbreak{}profiles/\allowbreak{}\_\allowbreak{}nemotron.env,\allowbreak{} ops/\allowbreak{}profiles/\allowbreak{}nemotron-\allowbreak{}super-\allowbreak{}vl-\allowbreak{}exl3-\allowbreak{}2.0.env,\allowbreak{} ops/\allowbreak{}profiles/\allowbreak{}nemotron-\allowbreak{}super-\allowbreak{}vl-\allowbreak{}exl3-\allowbreak{}4.0.env and 2 more} \\
\texttt{QSE\_\allowbreak{}NH\_\allowbreak{}GGUF} & \texttt{engine/\allowbreak{}models/\allowbreak{}nemotron\_\allowbreak{}h/\allowbreak{}load.py,\allowbreak{} ops/\allowbreak{}profiles/\allowbreak{}nemotron-\allowbreak{}bolt-\allowbreak{}iq2.env,\allowbreak{} server/\allowbreak{}diskcache.py} \\
\texttt{QSE\_\allowbreak{}NH\_\allowbreak{}HEAD} & \texttt{ops/\allowbreak{}profiles/\allowbreak{}\_\allowbreak{}nemotron.env,\allowbreak{} ops/\allowbreak{}profiles/\allowbreak{}nemotron-\allowbreak{}bolt-\allowbreak{}exl3-\allowbreak{}2.0.env,\allowbreak{} ops/\allowbreak{}profiles/\allowbreak{}nemotron-\allowbreak{}bolt-\allowbreak{}exl3-\allowbreak{}2.5.env and 10 more} \\
\texttt{QSE\_\allowbreak{}NH\_\allowbreak{}LILICORR} & \texttt{engine/\allowbreak{}models/\allowbreak{}nemotron\_\allowbreak{}h/\allowbreak{}serving.py,\allowbreak{} ops/\allowbreak{}profiles/\allowbreak{}nemotron-\allowbreak{}bolt-\allowbreak{}exl3-\allowbreak{}2.0.env,\allowbreak{} ops/\allowbreak{}profiles/\allowbreak{}nemotron-\allowbreak{}bolt-\allowbreak{}exl3-\allowbreak{}2.5.env and 8 more} \\
\texttt{QSE\_\allowbreak{}NH\_\allowbreak{}LOOKUP} & \texttt{engine/\allowbreak{}models/\allowbreak{}nemotron\_\allowbreak{}h/\allowbreak{}serving.py} \\
\texttt{QSE\_\allowbreak{}NH\_\allowbreak{}MODEL} & \texttt{engine/\allowbreak{}models/\allowbreak{}nemotron\_\allowbreak{}h/\allowbreak{}serving.py,\allowbreak{} ops/\allowbreak{}profiles/\allowbreak{}\_\allowbreak{}nemotron.env,\allowbreak{} ops/\allowbreak{}profiles/\allowbreak{}nemotron-\allowbreak{}bolt-\allowbreak{}exl3-\allowbreak{}2.0.env and 11 more} \\
\texttt{QSE\_\allowbreak{}NH\_\allowbreak{}VARIANT} & \texttt{ops/\allowbreak{}profiles/\allowbreak{}\_\allowbreak{}nemotron.env,\allowbreak{} ops/\allowbreak{}profiles/\allowbreak{}nemotron-\allowbreak{}bolt-\allowbreak{}exl3-\allowbreak{}2.0.env,\allowbreak{} ops/\allowbreak{}profiles/\allowbreak{}nemotron-\allowbreak{}bolt-\allowbreak{}exl3-\allowbreak{}2.5.env and 11 more} \\
\texttt{QSE\_\allowbreak{}NO\_\allowbreak{}ROLLBACK} & \texttt{ops/\allowbreak{}use-\allowbreak{}profile.sh,\allowbreak{} server/\allowbreak{}setup/\allowbreak{}apply.py} \\
\texttt{QSE\_\allowbreak{}OPS\_\allowbreak{}DIR} & \texttt{ops/\allowbreak{}start.sh,\allowbreak{} server/\allowbreak{}setup/\allowbreak{}apply.py,\allowbreak{} server/\allowbreak{}setup/\allowbreak{}device.py} \\
\texttt{QSE\_\allowbreak{}ORN\_\allowbreak{}DFLASH2} & \texttt{engine/\allowbreak{}models/\allowbreak{}qwen3\_\allowbreak{}5\_\allowbreak{}moe/\allowbreak{}serving.py,\allowbreak{} ops/\allowbreak{}profiles/\allowbreak{}\_\allowbreak{}ornith.env,\allowbreak{} ops/\allowbreak{}profiles/\allowbreak{}ornith-\allowbreak{}1.5-\allowbreak{}122b-\allowbreak{}exl3-\allowbreak{}4.0.env} \\
\texttt{QSE\_\allowbreak{}ORN\_\allowbreak{}DRAFTER} & \texttt{ops/\allowbreak{}profiles/\allowbreak{}\_\allowbreak{}ornith.env} \\
\texttt{QSE\_\allowbreak{}ORN\_\allowbreak{}HEAD} & \texttt{ops/\allowbreak{}profiles/\allowbreak{}\_\allowbreak{}ornith.env,\allowbreak{} ops/\allowbreak{}profiles/\allowbreak{}ornith-\allowbreak{}1.5-\allowbreak{}122b-\allowbreak{}exl3-\allowbreak{}4.0.env} \\
\texttt{QSE\_\allowbreak{}ORN\_\allowbreak{}VARIANT} & \texttt{ops/\allowbreak{}profiles/\allowbreak{}\_\allowbreak{}ornith.env,\allowbreak{} ops/\allowbreak{}profiles/\allowbreak{}ornith-\allowbreak{}1.5-\allowbreak{}122b-\allowbreak{}exl3-\allowbreak{}4.0.env} \\
\texttt{QSE\_\allowbreak{}PAGE\_\allowbreak{}CACHE\_\allowbreak{}DROP} & \texttt{ops/\allowbreak{}profiles/\allowbreak{}\_\allowbreak{}flashnext.env,\allowbreak{} ops/\allowbreak{}profiles/\allowbreak{}\_\allowbreak{}lfm25.env,\allowbreak{} ops/\allowbreak{}profiles/\allowbreak{}\_\allowbreak{}nemotron.env and 4 more} \\
\texttt{QSE\_\allowbreak{}PAGE\_\allowbreak{}CACHE\_\allowbreak{}KEEP} & \texttt{ops/\allowbreak{}profiles/\allowbreak{}\_\allowbreak{}flashnext.env,\allowbreak{} ops/\allowbreak{}serve.env} \\
\texttt{QSE\_\allowbreak{}PARALLEL} & \texttt{ops/\allowbreak{}profiles/\allowbreak{}\_\allowbreak{}flashnext.env,\allowbreak{} ops/\allowbreak{}profiles/\allowbreak{}\_\allowbreak{}lfm25.env,\allowbreak{} ops/\allowbreak{}profiles/\allowbreak{}\_\allowbreak{}nemotron.env and 4 more} \\
\texttt{QSE\_\allowbreak{}PARALLEL\_\allowbreak{}BATCH} & \texttt{server/\allowbreak{}app.py} \\
\texttt{QSE\_\allowbreak{}PARALLEL\_\allowbreak{}CUT} & \texttt{server/\allowbreak{}app.py} \\
\texttt{QSE\_\allowbreak{}PARALLEL\_\allowbreak{}DRAFT} & \texttt{server/\allowbreak{}app.py} \\
\texttt{QSE\_\allowbreak{}PARALLEL\_\allowbreak{}FLOOR\_\allowbreak{}GB} & \texttt{engine/\allowbreak{}kvpool.py,\allowbreak{} ops/\allowbreak{}profiles/\allowbreak{}\_\allowbreak{}flashnext.env,\allowbreak{} ops/\allowbreak{}serve.env and 1 more} \\
\texttt{QSE\_\allowbreak{}PARALLEL\_\allowbreak{}LAYER\_\allowbreak{}YIELD} & \texttt{server/\allowbreak{}app.py} \\
\texttt{QSE\_\allowbreak{}PARALLEL\_\allowbreak{}SWITCH\_\allowbreak{}MS} & \texttt{server/\allowbreak{}app.py} \\
\texttt{QSE\_\allowbreak{}PARALLEL\_\allowbreak{}TIMING} & \texttt{server/\allowbreak{}app.py} \\
\texttt{QSE\_\allowbreak{}PARALLEL\_\allowbreak{}WINDOW\_\allowbreak{}ROWS} & \texttt{server/\allowbreak{}app.py} \\
\texttt{QSE\_\allowbreak{}PRICE\_\allowbreak{}TABLE} & \texttt{server/\allowbreak{}app.py} \\
\texttt{QSE\_\allowbreak{}PROFILE\_\allowbreak{}NAME} & \texttt{ops/\allowbreak{}serve.env} \\
\texttt{QSE\_\allowbreak{}PROFILE\_\allowbreak{}OVERRIDE} & \texttt{ops/\allowbreak{}serve.env,\allowbreak{} ops/\allowbreak{}use-\allowbreak{}profile.sh,\allowbreak{} server/\allowbreak{}setup/\allowbreak{}apply.py and 1 more} \\
\texttt{QSE\_\allowbreak{}ROLES\_\allowbreak{}DYNAMIC} & \texttt{ops/\allowbreak{}serve.env,\allowbreak{} server/\allowbreak{}rolemem.py} \\
\texttt{QSE\_\allowbreak{}ROLE\_\allowbreak{}MAX\_\allowbreak{}HOLD\_\allowbreak{}S} & \texttt{ops/\allowbreak{}serve.env,\allowbreak{} server/\allowbreak{}rolemem.py} \\
\texttt{QSE\_\allowbreak{}ROLE\_\allowbreak{}RELOAD\_\allowbreak{}MARGIN\_\allowbreak{}GB} & \texttt{ops/\allowbreak{}serve.env,\allowbreak{} server/\allowbreak{}rolemem.py} \\
\texttt{QSE\_\allowbreak{}ROLE\_\allowbreak{}RELOAD\_\allowbreak{}S} & \texttt{ops/\allowbreak{}serve.env,\allowbreak{} server/\allowbreak{}rolemem.py} \\
\texttt{QSE\_\allowbreak{}ROLE\_\allowbreak{}THRASH\_\allowbreak{}S} & \texttt{ops/\allowbreak{}serve.env,\allowbreak{} server/\allowbreak{}rolemem.py} \\
\texttt{QSE\_\allowbreak{}ROLE\_\allowbreak{}TICK\_\allowbreak{}FAKE} & \texttt{server/\allowbreak{}app.py} \\
\texttt{QSE\_\allowbreak{}ROLE\_\allowbreak{}TICK\_\allowbreak{}S} & \texttt{server/\allowbreak{}rolemem.py} \\
\texttt{QSE\_\allowbreak{}ROLE\_\allowbreak{}UNLOAD\_\allowbreak{}LAYA} & \texttt{ops/\allowbreak{}serve.env,\allowbreak{} server/\allowbreak{}app.py,\allowbreak{} server/\allowbreak{}rolemem.py} \\
\texttt{QSE\_\allowbreak{}SANDBOX\_\allowbreak{}SOCKET} & \texttt{ops/\allowbreak{}sandbox/\allowbreak{}install.sh} \\
\texttt{QSE\_\allowbreak{}SECRETS} & \texttt{install.sh,\allowbreak{} ops/\allowbreak{}chatd.sh,\allowbreak{} ops/\allowbreak{}make-\allowbreak{}secrets.sh and 4 more} \\
\texttt{QSE\_\allowbreak{}SESSION\_\allowbreak{}S} & \texttt{server/\allowbreak{}auth.py} \\
\texttt{QSE\_\allowbreak{}SETUP\_\allowbreak{}ENV} & \texttt{ops/\allowbreak{}serve.env,\allowbreak{} server/\allowbreak{}setup/\allowbreak{}apply.py} \\
\texttt{QSE\_\allowbreak{}SETUP\_\allowbreak{}ENV\_\allowbreak{}FILE} & \texttt{ops/\allowbreak{}serve.env} \\
\texttt{QSE\_\allowbreak{}SETUP\_\allowbreak{}MODE} & \texttt{ops/\allowbreak{}start.sh} \\
\texttt{QSE\_\allowbreak{}SIDE\_\allowbreak{}GRAPHS} & \texttt{ops/\allowbreak{}serve.env,\allowbreak{} server/\allowbreak{}app.py} \\
\texttt{QSE\_\allowbreak{}SIDE\_\allowbreak{}MAX\_\allowbreak{}LEN} & \texttt{ops/\allowbreak{}profiles/\allowbreak{}\_\allowbreak{}flashnext.env,\allowbreak{} ops/\allowbreak{}serve.env,\allowbreak{} server/\allowbreak{}app.py} \\
\texttt{QSE\_\allowbreak{}SIDE\_\allowbreak{}RESIDENT\_\allowbreak{}GB} & \texttt{server/\allowbreak{}app.py} \\
\texttt{QSE\_\allowbreak{}SIDE\_\allowbreak{}SLOT\_\allowbreak{}CONST\_\allowbreak{}GB} & \texttt{ops/\allowbreak{}serve.env,\allowbreak{} server/\allowbreak{}app.py} \\
\texttt{QSE\_\allowbreak{}STATE\_\allowbreak{}DIR} & \texttt{ops/\allowbreak{}chatd.sh,\allowbreak{} ops/\allowbreak{}make-\allowbreak{}secrets.sh,\allowbreak{} ops/\allowbreak{}start.sh and 3 more} \\
\texttt{QSE\_\allowbreak{}STATUS\_\allowbreak{}FILE} & \texttt{ops/\allowbreak{}engine\_\allowbreak{}status.py} \\
\texttt{QSE\_\allowbreak{}STATUS\_\allowbreak{}HOST} & \texttt{ops/\allowbreak{}engine\_\allowbreak{}status.py,\allowbreak{} ops/\allowbreak{}hold.sh,\allowbreak{} ops/\allowbreak{}serve.env and 1 more} \\
\texttt{QSE\_\allowbreak{}STATUS\_\allowbreak{}KEEP} & \texttt{ops/\allowbreak{}hold.sh,\allowbreak{} ops/\allowbreak{}stop.sh} \\
\texttt{QSE\_\allowbreak{}STATUS\_\allowbreak{}PORT} & \texttt{install.sh,\allowbreak{} ops/\allowbreak{}engine\_\allowbreak{}status.py,\allowbreak{} ops/\allowbreak{}hold.sh and 4 more} \\
\texttt{QSE\_\allowbreak{}STATUS\_\allowbreak{}TREE} & \texttt{install.sh,\allowbreak{} ops/\allowbreak{}engine\_\allowbreak{}status.py,\allowbreak{} ops/\allowbreak{}hold.sh} \\
\texttt{QSE\_\allowbreak{}TEST} & \texttt{proxy/\allowbreak{}app.py,\allowbreak{} server/\allowbreak{}app.py,\allowbreak{} server/\allowbreak{}ledger.py} \\
\texttt{QSE\_\allowbreak{}USAGE\_\allowbreak{}DEFAULT} & \texttt{server/\allowbreak{}app.py} \\
\texttt{QSE\_\allowbreak{}USAGE\_\allowbreak{}LEDGER} & \texttt{install.sh,\allowbreak{} ops/\allowbreak{}serve.env,\allowbreak{} ops/\allowbreak{}start.sh and 3 more} \\
\texttt{QSE\_\allowbreak{}WEIGHTS} & \texttt{engine/\allowbreak{}models/\allowbreak{}flashnext/\allowbreak{}serving.py,\allowbreak{} engine/\allowbreak{}models/\allowbreak{}lfm2\_\allowbreak{}moe/\allowbreak{}serving.py,\allowbreak{} engine/\allowbreak{}models/\allowbreak{}nemotron\_\allowbreak{}h/\allowbreak{}serving.py and 15 more} \\
\texttt{QSE\_\allowbreak{}WEIGHTS\_\allowbreak{}ALSO} & \texttt{ops/\allowbreak{}profiles/\allowbreak{}\_\allowbreak{}flashnext.env,\allowbreak{} ops/\allowbreak{}use-\allowbreak{}profile.sh} \\
\texttt{QSE\_\allowbreak{}WEIGHTS\_\allowbreak{}DIR} & \texttt{ops/\allowbreak{}profiles/\allowbreak{}\_\allowbreak{}exl3.env,\allowbreak{} ops/\allowbreak{}profiles/\allowbreak{}\_\allowbreak{}flashnext.env,\allowbreak{} ops/\allowbreak{}profiles/\allowbreak{}\_\allowbreak{}lfm25.env and 9 more} \\
\texttt{QSE\_\allowbreak{}WEIGHTS\_\allowbreak{}FILES} & \texttt{ops/\allowbreak{}profiles/\allowbreak{}\_\allowbreak{}exl3.env,\allowbreak{} ops/\allowbreak{}profiles/\allowbreak{}\_\allowbreak{}flashnext.env,\allowbreak{} ops/\allowbreak{}profiles/\allowbreak{}\_\allowbreak{}lfm25.env and 7 more} \\
\texttt{QSE\_\allowbreak{}WEIGHTS\_\allowbreak{}GB} & \texttt{ops/\allowbreak{}profiles/\allowbreak{}exl3-\allowbreak{}2.5.env,\allowbreak{} ops/\allowbreak{}profiles/\allowbreak{}exl3-\allowbreak{}3.0.env,\allowbreak{} ops/\allowbreak{}profiles/\allowbreak{}exl3-\allowbreak{}3.5.env and 31 more} \\
\texttt{QSE\_\allowbreak{}WEIGHTS\_\allowbreak{}HF} & \texttt{ops/\allowbreak{}profiles/\allowbreak{}\_\allowbreak{}exl3.env,\allowbreak{} ops/\allowbreak{}profiles/\allowbreak{}\_\allowbreak{}flashnext.env,\allowbreak{} ops/\allowbreak{}profiles/\allowbreak{}lfm25-\allowbreak{}bf16.env and 23 more} \\
\texttt{QSE\_\allowbreak{}WEIGHTS\_\allowbreak{}REVISION} & \texttt{ops/\allowbreak{}profiles/\allowbreak{}\_\allowbreak{}flashnext.env,\allowbreak{} ops/\allowbreak{}profiles/\allowbreak{}flashnext-\allowbreak{}2.05.env,\allowbreak{} ops/\allowbreak{}profiles/\allowbreak{}flashnext-\allowbreak{}3.05.env and 26 more} \\
\texttt{QSE\_\allowbreak{}WREN\_\allowbreak{}MODEL} & \texttt{server/\allowbreak{}setup/\allowbreak{}wren.py} \\
\end{reftabletwo}


The most used: \code{QSE_CONDUCTOR_CONFIG}\index{qseconductorconfig@\texttt{QSE\_CONDUCTOR\_CONFIG}}, \code{QSE_USAGE_LEDGER}\index{qseusageledger@\texttt{QSE\_USAGE\_LEDGER}}, \code{QSE_DISK_CACHE}\index{qsediskcache@\texttt{QSE\_DISK\_CACHE}} and \code{QSE_BACKENDS}\index{qsebackends@\texttt{QSE\_BACKENDS}} name the conductor's directory, the ledger, the disk tier's directory and the proxy's config database; \code{ops/start.sh} passes them to the service only, so a benchmark server started with \code{serve.env} exported never writes them. \code{QSE_DASHBOARD_LOGIN}\index{qsedashboardlogin@\texttt{QSE\_DASHBOARD\_LOGIN}} turns the dashboard login on. \code{QSE_ADMIN_TOKEN} and \code{QSE_METRICS_TOKEN} live in \code{secrets.env}, which \code{ops/make-secrets.sh} writes.

\section{Proxy, collector and plugin variables}\index{TANDEM variables}

% Generated by manual/tools/gen_reference.py from the code. Do not edit by hand.

\begin{reftable}{TANDEM\_* variables: the proxy, the collector, the plugins and the installer}{tab:ref-tandem}{\hd{Variable} & \hd{Example} & \hd{What it sets}}
\texttt{TANDEM\_\allowbreak{}BACKENDS} & \texttt{deepseek,\allowbreak{}vllm1} & the backends the environment declares \\
\texttt{TANDEM\_\allowbreak{}BACKEND\_\allowbreak{}\textless{}ID\textgreater{}\_\allowbreak{}KIND} & \texttt{deepseek} & generic \textbar{} tandemllm \textbar{} vllm \textbar{} sglang \textbar{} openai \textbar{} deepseek \\
\texttt{TANDEM\_\allowbreak{}BACKEND\_\allowbreak{}\textless{}ID\textgreater{}\_\allowbreak{}URL} & \texttt{https:\allowbreak{}/\allowbreak{}/\allowbreak{}api.deepseek.com/\allowbreak{}v1} &  \\
\texttt{TANDEM\_\allowbreak{}BACKEND\_\allowbreak{}\textless{}ID\textgreater{}\_\allowbreak{}KEY} & \texttt{(a key)} & held in memory, never stored \\
\texttt{TANDEM\_\allowbreak{}BACKEND\_\allowbreak{}\textless{}ID\textgreater{}\_\allowbreak{}LABEL} & \texttt{DeepSeek} &  \\
\texttt{TANDEM\_\allowbreak{}BACKEND\_\allowbreak{}\textless{}ID\textgreater{}\_\allowbreak{}MODELS} & \texttt{deepseek-\allowbreak{}flash,\allowbreak{}deepseek-\allowbreak{}v4-\allowbreak{}pro} &  \\
\texttt{TANDEM\_\allowbreak{}BACKEND\_\allowbreak{}\textless{}ID\textgreater{}\_\allowbreak{}MODELS\_\allowbreak{}MODE} & \texttt{discover} & discover (GET /models, plus MODELS) \textbar{} list \\
\texttt{TANDEM\_\allowbreak{}BACKEND\_\allowbreak{}\textless{}ID\textgreater{}\_\allowbreak{}FALLBACK} & \texttt{openrouter} & backends tried after this one, in order \\
\texttt{TANDEM\_\allowbreak{}BACKEND\_\allowbreak{}\textless{}ID\textgreater{}\_\allowbreak{}MAX\_\allowbreak{}INFLIGHT} & \texttt{none} & requests at once, past it a fallback or a 503; none (the default): no limit from the proxy, the backend takes what it accepts \\
\texttt{TANDEM\_\allowbreak{}BACKEND\_\allowbreak{}\textless{}ID\textgreater{}\_\allowbreak{}ENABLED} & \texttt{1} &  \\
\texttt{TANDEM\_\allowbreak{}BACKEND\_\allowbreak{}\textless{}ID\textgreater{}\_\allowbreak{}INJECT\_\allowbreak{}USAGE} & \texttt{1} & ask a stream for its usage frame \\
\texttt{TANDEM\_\allowbreak{}BACKEND\_\allowbreak{}\textless{}ID\textgreater{}\_\allowbreak{}HARNESS} & \texttt{off} & off \textbar{} shadow \textbar{} on \\
\texttt{TANDEM\_\allowbreak{}BACKEND\_\allowbreak{}\textless{}ID\textgreater{}\_\allowbreak{}HARNESS\_\allowbreak{}PROFILE} & \texttt{coding} & the conductor profile of its requests \\
\texttt{TANDEM\_\allowbreak{}BACKEND\_\allowbreak{}\textless{}ID\textgreater{}\_\allowbreak{}WRITER} & \texttt{auto} & the background memory writer on this backend (proxy/memwriter.py): auto (on, off for a paid cloud model) \textbar{} off \textbar{} shadow \textbar{} on \\
\texttt{TANDEM\_\allowbreak{}BACKEND\_\allowbreak{}\textless{}ID\textgreater{}\_\allowbreak{}WRITER\_\allowbreak{}USD\_\allowbreak{}DAY} & \texttt{0.25} & its spend cap a day on a paid model \\
\texttt{TANDEM\_\allowbreak{}BACKEND\_\allowbreak{}\textless{}ID\textgreater{}\_\allowbreak{}WRITER\_\allowbreak{}RUNS\_\allowbreak{}DAY} & \texttt{48} & its runs a day \\
\texttt{TANDEM\_\allowbreak{}BACKEND\_\allowbreak{}\textless{}ID\textgreater{}\_\allowbreak{}WRITER\_\allowbreak{}PROMPT\_\allowbreak{}TOKENS} & \texttt{32000} & its prompt at most (older messages dropped) \\
\texttt{TANDEM\_\allowbreak{}BACKEND\_\allowbreak{}\textless{}ID\textgreater{}\_\allowbreak{}TIMEOUT\_\allowbreak{}CONNECT} & \texttt{10} & seconds; also \_TIMEOUT\_FIRST\_BYTE, \_TIMEOUT\_IDLE \\
\texttt{TANDEM\_\allowbreak{}BACKEND\_\allowbreak{}\textless{}ID\textgreater{}\_\allowbreak{}HOST} & \texttt{auto} & where the System view reads the backend's machine: auto (this machine when the URL is), same, collector, exporter, off \\
\texttt{TANDEM\_\allowbreak{}BACKEND\_\allowbreak{}\textless{}ID\textgreater{}\_\allowbreak{}HOST\_\allowbreak{}URL} & \texttt{(a URL)} & the remote collector (python -m proxy.hostcollect), or a node\_exporter page \\
\texttt{TANDEM\_\allowbreak{}BACKEND\_\allowbreak{}\textless{}ID\textgreater{}\_\allowbreak{}HOST\_\allowbreak{}DCGM\_\allowbreak{}URL} & \texttt{(a URL)} & a DCGM exporter page (mode exporter) \\
\texttt{TANDEM\_\allowbreak{}COLLECTOR\_\allowbreak{}TOKEN} & \texttt{...} & the remote collectors' bearer token, in memory only (proxy/hostcollect.py) \\
\texttt{TANDEM\_\allowbreak{}ROUTE\_\allowbreak{}\textless{}n\textgreater{}} & \texttt{gpt-\allowbreak{}*=\allowbreak{}openai} & routing rules, ordered by n \\
\texttt{TANDEM\_\allowbreak{}PRICE\_\allowbreak{}\textless{}n\textgreater{}} & \texttt{deepseek/\allowbreak{}deepseek-\allowbreak{}flash=\allowbreak{}0.30,\allowbreak{}0.006,\allowbreak{}1.2} & USD per million: input, cached input, output, and optionally cache write \\
\texttt{TANDEM\_\allowbreak{}PRICE\_\allowbreak{}\textless{}n\textgreater{}\_\allowbreak{}OFFPEAK} & \texttt{0.15,\allowbreak{}0.003,\allowbreak{}0.6} & the same model outside the provider's peak hours (DeepSeek) \\
\texttt{TANDEM\_\allowbreak{}API\_\allowbreak{}KEYS} & \texttt{(a key)} & the proxy's own keys, in memory only \\
\texttt{TANDEM\_\allowbreak{}SECRETS\_\allowbreak{}KEY} & \texttt{\textless{}base64} & of 32 bytes\textgreater{} the master key (proxy/secrets.py) \\
\texttt{TANDEM\_\allowbreak{}ADMIN\_\allowbreak{}TOKEN} & & read in \texttt{ops/\allowbreak{}dash.env.example,\allowbreak{} proxy/\allowbreak{}\_\allowbreak{}\_\allowbreak{}main\_\allowbreak{}\_\allowbreak{}.py} \\
\texttt{TANDEM\_\allowbreak{}BACKEND\_\allowbreak{}DEEPSEEK\_\allowbreak{}KEY} & & read in \texttt{ops/\allowbreak{}dash.env.example} \\
\texttt{TANDEM\_\allowbreak{}BACKEND\_\allowbreak{}DEEPSEEK\_\allowbreak{}KIND} & & read in \texttt{ops/\allowbreak{}dash.env.example} \\
\texttt{TANDEM\_\allowbreak{}BACKEND\_\allowbreak{}DEEPSEEK\_\allowbreak{}LABEL} & & read in \texttt{ops/\allowbreak{}dash.env.example} \\
\texttt{TANDEM\_\allowbreak{}BACKEND\_\allowbreak{}VLLM1\_\allowbreak{}FALLBACK} & & read in \texttt{ops/\allowbreak{}dash.env.example} \\
\texttt{TANDEM\_\allowbreak{}BACKEND\_\allowbreak{}VLLM1\_\allowbreak{}HOST} & & read in \texttt{ops/\allowbreak{}dash.env.example} \\
\texttt{TANDEM\_\allowbreak{}BACKEND\_\allowbreak{}VLLM1\_\allowbreak{}HOST\_\allowbreak{}DCGM\_\allowbreak{}URL} & & read in \texttt{ops/\allowbreak{}dash.env.example} \\
\texttt{TANDEM\_\allowbreak{}BACKEND\_\allowbreak{}VLLM1\_\allowbreak{}HOST\_\allowbreak{}URL} & & read in \texttt{ops/\allowbreak{}dash.env.example} \\
\texttt{TANDEM\_\allowbreak{}BACKEND\_\allowbreak{}VLLM1\_\allowbreak{}KIND} & & read in \texttt{ops/\allowbreak{}dash.env.example} \\
\texttt{TANDEM\_\allowbreak{}BACKEND\_\allowbreak{}VLLM1\_\allowbreak{}MAX\_\allowbreak{}INFLIGHT} & & read in \texttt{ops/\allowbreak{}dash.env.example} \\
\texttt{TANDEM\_\allowbreak{}BACKEND\_\allowbreak{}VLLM1\_\allowbreak{}TIMEOUT\_\allowbreak{}FIRST\_\allowbreak{}BYTE} & & read in \texttt{ops/\allowbreak{}dash.env.example} \\
\texttt{TANDEM\_\allowbreak{}BACKEND\_\allowbreak{}VLLM1\_\allowbreak{}URL} & & read in \texttt{ops/\allowbreak{}dash.env.example} \\
\texttt{TANDEM\_\allowbreak{}COLLECT} & & read in \texttt{proxy/\allowbreak{}launch.py} \\
\texttt{TANDEM\_\allowbreak{}COLLECT\_\allowbreak{}BASE} & & read in \texttt{proxy/\allowbreak{}launch.py} \\
\texttt{TANDEM\_\allowbreak{}COLLECT\_\allowbreak{}TOKEN} & & read in \texttt{plugins/\allowbreak{}sglang/\allowbreak{}tandem\_\allowbreak{}sglang\_\allowbreak{}plugin/\allowbreak{}\_\allowbreak{}push.py,\allowbreak{} plugins/\allowbreak{}vllm/\allowbreak{}tandem\_\allowbreak{}vllm\_\allowbreak{}plugin/\allowbreak{}\_\allowbreak{}\_\allowbreak{}init\_\allowbreak{}\_\allowbreak{}.py} \\
\texttt{TANDEM\_\allowbreak{}COLLECT\_\allowbreak{}URL} & & read in \texttt{plugins/\allowbreak{}sglang/\allowbreak{}tandem\_\allowbreak{}sglang\_\allowbreak{}plugin/\allowbreak{}\_\allowbreak{}push.py,\allowbreak{} plugins/\allowbreak{}vllm/\allowbreak{}tandem\_\allowbreak{}vllm\_\allowbreak{}plugin/\allowbreak{}\_\allowbreak{}\_\allowbreak{}init\_\allowbreak{}\_\allowbreak{}.py} \\
\texttt{TANDEM\_\allowbreak{}CONDUCTOR} & & read in \texttt{proxy/\allowbreak{}\_\allowbreak{}\_\allowbreak{}main\_\allowbreak{}\_\allowbreak{}.py} \\
\texttt{TANDEM\_\allowbreak{}CONFIG\_\allowbreak{}DB} & & read in \texttt{ops/\allowbreak{}dash.env.example,\allowbreak{} proxy/\allowbreak{}\_\allowbreak{}\_\allowbreak{}main\_\allowbreak{}\_\allowbreak{}.py} \\
\texttt{TANDEM\_\allowbreak{}DASHBOARD\_\allowbreak{}LOGIN} & & read in \texttt{ops/\allowbreak{}dash.env.example,\allowbreak{} proxy/\allowbreak{}\_\allowbreak{}\_\allowbreak{}main\_\allowbreak{}\_\allowbreak{}.py} \\
\texttt{TANDEM\_\allowbreak{}DIR} & & read in \texttt{install.sh} \\
\texttt{TANDEM\_\allowbreak{}HEALTH\_\allowbreak{}TIMEOUT} & & read in \texttt{install.sh} \\
\texttt{TANDEM\_\allowbreak{}HOST} & & read in \texttt{install.sh,\allowbreak{} ops/\allowbreak{}dash.env.example} \\
\texttt{TANDEM\_\allowbreak{}IDLE\_\allowbreak{}TIMEOUT\_\allowbreak{}S} & & read in \texttt{proxy/\allowbreak{}\_\allowbreak{}\_\allowbreak{}main\_\allowbreak{}\_\allowbreak{}.py,\allowbreak{} proxy/\allowbreak{}app.py} \\
\texttt{TANDEM\_\allowbreak{}MIN\_\allowbreak{}FREE\_\allowbreak{}GIB} & & read in \texttt{install.sh} \\
\texttt{TANDEM\_\allowbreak{}MODEL\_\allowbreak{}NAME} & & read in \texttt{install.sh} \\
\texttt{TANDEM\_\allowbreak{}NOTIFY\_\allowbreak{}URL} & & read in \texttt{server/\allowbreak{}conductor/\allowbreak{}monitor/\allowbreak{}notify.py,\allowbreak{} server/\allowbreak{}conductor/\allowbreak{}monitor/\allowbreak{}sessiontasks.py} \\
\texttt{TANDEM\_\allowbreak{}PORT} & & read in \texttt{install.sh,\allowbreak{} ops/\allowbreak{}dash.env.example} \\
\texttt{TANDEM\_\allowbreak{}PRICE\_\allowbreak{}1} & & read in \texttt{ops/\allowbreak{}dash.env.example} \\
\texttt{TANDEM\_\allowbreak{}PRICE\_\allowbreak{}1\_\allowbreak{}OFFPEAK} & & read in \texttt{ops/\allowbreak{}dash.env.example} \\
\texttt{TANDEM\_\allowbreak{}REF} & & read in \texttt{install.sh} \\
\texttt{TANDEM\_\allowbreak{}REPO} & & read in \texttt{install.sh} \\
\texttt{TANDEM\_\allowbreak{}ROUTE\_\allowbreak{}1} & & read in \texttt{ops/\allowbreak{}dash.env.example} \\
\texttt{TANDEM\_\allowbreak{}SCRATCH} & & read in \texttt{server/\allowbreak{}conductor/\allowbreak{}packages.py} \\
\texttt{TANDEM\_\allowbreak{}SKIP\_\allowbreak{}MEMCHECK} & & read in \texttt{install.sh} \\
\texttt{TANDEM\_\allowbreak{}URL} & & read in \texttt{server/\allowbreak{}conductor/\allowbreak{}packages.py,\allowbreak{} server/\allowbreak{}laya\_\allowbreak{}kit/\allowbreak{}laya\_\allowbreak{}drive.py} \\
\texttt{TANDEM\_\allowbreak{}USAGE\_\allowbreak{}LEDGER} & & read in \texttt{ops/\allowbreak{}dash.env.example,\allowbreak{} proxy/\allowbreak{}\_\allowbreak{}\_\allowbreak{}main\_\allowbreak{}\_\allowbreak{}.py} \\
\end{reftable}


\section{The conductor's documents}\index{conductor documents}

\begin{table}[htbp]
\centering\small
\caption{The conductor's versioned documents.}\label{tab:ref-documents}
\begin{tabularx}{\linewidth}{@{}L{.16\linewidth}YL{.22\linewidth}@{}}
\toprule
\hd{document} & \hd{what it holds} & \hd{edited on} \\
\midrule
\code{profiles} & one document per profile: mode, decisions, task modes and caps, thinking, \code{ui}, outcomes, library, memory, shadow & Profiles \\
\code{router} & the classes and their policy, the classifier stages and thresholds, the structural rules, side shapes, outcomes, library & Routing \\
\code{packages} & every package: source, switch, limits, flags & Packages \\
\code{roles} & every role's contract & Profiles, Roles \\
\code{settings} & the default profile, transcripts, the disk tier, memory notes and the writer & Memory, Profiles \\
\code{deciders} & the decision models, their questions, Laya for scripts & Deciders \\
\code{residency} & the floor, the margin, the settle time, the load order, the claims & Profiles, Residency \\
\code{monitor} & the session monitor's policy (version 0 is \code{ops/supervisor.yaml}) & Monitor \\
\code{harness} & the harness's settings per backend, behind the proxy & Backends \\
\bottomrule
\end{tabularx}
\end{table}

A document is JSON, checked against the conductor's schema (\code{GET /v1/conductor/meta} serves it, and \code{docs/contract/conductor-v1/} holds a copy). A version stored by an older release loads with defaults for every field it lacks, and the journal says so once.
